%==============================================================
%  INTRODUCCIÓN A SQL Y BASES DE DATOS RELACIONALES
%  Módulo 1 – Fundamentos teóricos
%  Formato: Artículo científico / Material académico
%==============================================================

\documentclass[12pt, a4paper]{article}

% ─── Codificación y lenguaje ──────────────────────────────
\usepackage[utf8]{inputenc}
\usepackage[T1]{fontenc}
\usepackage[spanish, es-tabla]{babel}

% ─── Tipografía ───────────────────────────────────────────
\usepackage{lmodern}
\usepackage{microtype}          % Mejora el espaciado tipográfico

% ─── Matemáticas y símbolos ───────────────────────────────
\usepackage{amsmath, amssymb}

% ─── Geometría de página ──────────────────────────────────
\usepackage[
  top=2.5cm, bottom=2.5cm,
  left=3cm, right=3cm
]{geometry}

% ─── Gráficos y figuras ───────────────────────────────────
\usepackage{graphicx}
\usepackage{float}
\usepackage{caption}
\usepackage{subcaption}
\usepackage[dvipsnames, table]{xcolor}

% ─── Tablas ───────────────────────────────────────────────
\usepackage{booktabs}
\usepackage{array}
\usepackage{tabularx}
\usepackage{longtable}
\usepackage{multirow}

% ─── Listados de código ───────────────────────────────────
\usepackage{listings}
\usepackage{lstautogobble}

% Estilo general para código SQL
\lstdefinestyle{sql}{
  language=SQL,
  basicstyle=\ttfamily\small,
  keywordstyle=\bfseries\color{NavyBlue},
  stringstyle=\color{BrickRed},
  commentstyle=\itshape\color{Gray},
  numberstyle=\tiny\color{Gray},
  numbers=left,
  numbersep=8pt,
  stepnumber=1,
  backgroundcolor=\color{gray!6},
  frame=single,
  framerule=0.4pt,
  rulecolor=\color{gray!50},
  tabsize=4,
  showstringspaces=false,
  breaklines=true,
  breakatwhitespace=false,
  captionpos=b,
  autogobble=true,
  xleftmargin=0.5cm,
  morekeywords={
    DATABASE, DATABASES, TABLES, SHOW, USE, ENGINE, CHARSET,
    AUTO_INCREMENT, UNSIGNED, TINYINT, SMALLINT, MEDIUMINT,
    BIGINT, DECIMAL, FLOAT, DOUBLE, DATETIME, TIMESTAMP,
    VARCHAR, TEXT, BLOB, ENUM, SET, DESCRIBE, DESC,
    FOREIGN KEY, REFERENCES, ON DELETE, ON UPDATE,
    CASCADE, SET NULL, RESTRICT
  }
}
\lstset{style=sql}

% ─── Hipervínculos y referencias ──────────────────────────
\usepackage[
  colorlinks=true,
  linkcolor=NavyBlue,
  citecolor=NavyBlue,
  urlcolor=NavyBlue
]{hyperref}
\usepackage{cleveref}

% ─── Cabecera y pie de página ─────────────────────────────
\usepackage{fancyhdr}
\pagestyle{fancy}
\fancyhf{}
\fancyhead[L]{\small\textit{Comisión Nacional Bancaria y de Valores}}
\fancyhead[R]{\small\textit{Miyoshi Mikami }}
\fancyfoot[C]{\small\thepage}
\renewcommand{\headrulewidth}{0.4pt}

% ─── Cajas informativas ───────────────────────────────────
\usepackage{tcolorbox}
\tcbuselibrary{breakable, skins}

\newtcolorbox{definicion}[1]{
  colback=NavyBlue!5,
  colframe=NavyBlue!60,
  fonttitle=\bfseries\small,
  title=Definición: #1,
  breakable
}

\newtcolorbox{ejemplo}{
  colback=ForestGreen!5,
  colframe=ForestGreen!50,
  fonttitle=\bfseries\small,
  title=Ejemplo práctico,
  breakable
}

\newtcolorbox{nota}{
  colback=Goldenrod!10,
  colframe=Goldenrod!70,
  fonttitle=\bfseries\small,
  title=Nota importante,
  breakable
}

% ─── Enumeraciones y listas ───────────────────────────────
\usepackage{enumitem}
\setlist[itemize]{leftmargin=*, itemsep=2pt, topsep=4pt}
\setlist[enumerate]{leftmargin=*, itemsep=2pt, topsep=4pt}

% ─── Espaciado ────────────────────────────────────────────
\setlength{\parskip}{4pt}
\setlength{\parindent}{0pt}

% ─── Bibliografía ─────────────────────────────────────────
\usepackage[
  backend=biber,
  style=ieee,
  sorting=nyt
]{biblatex}
\addbibresource{referencias.bib}

%==============================================================
%  DATOS DEL DOCUMENTO
%==============================================================
\title{%
  {\large \textsc{Curso de SQL y Bases de Datos Relacionales}}\\[6pt]
  {\LARGE \textbf{SQL}}\\[4pt]
  {\large Inserción, Consulta y Adminitración de información en base de datos relacionales basadas en SQL}
}

\author{%
  \textbf{Material de Referencia }\\[2pt]
  \small \texttt{miyoshikamium@gmail.mx}
}

\date{\small Versión 4.5 — \today}

%==============================================================
%  INICIO DEL DOCUMENTO
%==============================================================
\begin{document}

\maketitle
\thispagestyle{fancy}

%--------------------------------------------------------------
\begin{abstract}
\noindent
Este documento constituye el material teórico de referencia para el
Módulo 1 del \textit{Curso de SQL y Bases de Datos Relacionales}.
Se abordan los fundamentos históricos y conceptuales que dieron origen
a los sistemas de gestión de bases de datos (SGBD) relacionales, la
evolución del lenguaje SQL desde sus orígenes en los laboratorios de
IBM hasta los estándares modernos, y los principios del modelo
relacional propuesto por Edgar F.\ Codd.
El objetivo es proveer al estudiante una base conceptual sólida que
permita contextualizar cada construcción sintáctica de SQL dentro de
un marco teórico coherente, facilitando un aprendizaje significativo
y duradero.

\smallskip
\textbf{Palabras clave:} base de datos, modelo relacional, SQL, SGBD,
álgebra relacional, normalización, ACID.
\end{abstract}

\tableofcontents
\vspace{8pt}
\hrule
\vspace{10pt}
%==============================================================
\section{El Problema de los Datos: Una Perspectiva Histórica}
%==============================================================

Desde los albores de la civilización, la humanidad ha necesitado
almacenar y recuperar información de manera sistemática.
Las tablillas de arcilla sumerias (~3000 a.C.), los registros
contables del Imperio Romano, o los libros de contabilidad de los
comerciantes medievales son, en esencia, \textit{sistemas de gestión
de información}: colecciones de datos organizados con un propósito
definido y un mecanismo de consulta.

La revolución informática de mediados del siglo XX transformó este
problema milenario en una disciplina formal.
Con la proliferación de computadoras en organizaciones bancarias,
gubernamentales y empresariales durante los años 1950 y 1960,
la necesidad de almacenar grandes volúmenes de datos de forma
eficiente y recuperarlos con rapidez se volvió crítica.

\subsection{Era pre-relacional (1950–1970)}

Los primeros sistemas de gestión de datos surgieron como soluciones
ad-hoc a problemas específicos.
Predominaban dos modelos principales:

\begin{itemize}
  \item \textbf{Modelo jerárquico:} Los datos se organizaban en una
    estructura de árbol. IBM desarrolló IMS
    (\textit{Information Management System}) en 1966 para la misión
    Apolo de la NASA~\cite{date2003introduction}.
    Aunque eficiente para consultas predefinidas, era rígido:
    agregar una nueva relación entre datos requería reestructurar
    físicamente el sistema.

  \item \textbf{Modelo de red (CODASYL):} Propuesto por el comité
    CODASYL en 1969, permitía relaciones muchos-a-muchos mediante
    punteros físicos entre registros~\cite{codasyl1971}.
    Era más flexible que el jerárquico, pero extremadamente complejo
    de administrar y altamente dependiente de la estructura física
    del almacenamiento.
\end{itemize}

El problema fundamental de ambos modelos era la
\textbf{dependencia físico-lógica}: los programas de aplicación
debían conocer cómo estaban organizados físicamente los datos en
disco para poder acceder a ellos. Un cambio en la estructura de
almacenamiento implicaba reescribir todos los programas.

%==============================================================
\section{Edgar F.\ Codd y el Modelo Relacional}
%==============================================================

En junio de 1970, Edgar Frank Codd —matemático e investigador
de IBM— publicó el artículo que cambiaría para siempre el mundo de
las bases de datos:
\textit{``A Relational Model of Data for Large Shared Data Banks''}
en la revista \textit{Communications of the ACM}~\cite{codd1970}.

\begin{definicion}{Modelo Relacional Codd, 1970}
Un modelo de datos en el que la información se representa
exclusivamente mediante \textbf{relaciones} (tablas bidimensionales),
y en el que todas las operaciones sobre los datos se expresan
mediante un conjunto cerrado de operaciones del
\textbf{álgebra relacional}, sin necesidad de conocer la
organización física del almacenamiento.
\end{definicion}

La contribución de Codd fue revolucionaria por tres razones:

\begin{enumerate}
  \item \textbf{Independencia de datos:} Los programas de aplicación
    operan sobre una representación lógica (tablas), completamente
    desacoplada de cómo los datos se almacenan físicamente.

  \item \textbf{Fundamento matemático:} El modelo se basa en la
    teoría de conjuntos y la lógica de predicados de primer orden,
    lo que permite razonar formalmente sobre las operaciones.

  \item \textbf{Lenguaje declarativo:} El programador especifica
    \emph{qué} datos quiere, no \emph{cómo} obtenerlos.
    El sistema optimiza el acceso automáticamente.
\end{enumerate}

\subsection{Las 12 reglas de Codd}

En 1985, Codd formalizó su modelo en 12 reglas
(numeradas del 0 al 12) que debe cumplir un sistema para
ser considerado verdaderamente relacional~\cite{codd1985}.
Las más relevantes para este curso son:

\begin{table}[H]
\centering
\small
\renewcommand{\arraystretch}{1.3}
\begin{tabularx}{\columnwidth}{@{} c X @{}}
\toprule
\textbf{Regla} & \textbf{Descripción} \\
\midrule
0 & Gestión mediante capacidades relacionales \\
1 & Representación de la información como tablas \\
2 & Acceso garantizado a cada dato (tabla + clave + columna) \\
3 & Tratamiento sistemático de valores nulos \\
6 & Actualización de vistas \\
9 & Independencia de datos física y lógica \\
12 & No subversión del modelo relacional \\
\bottomrule
\end{tabularx}
\caption{Selección de las 12 reglas de Codd (1985).}
\label{tab:reglas-codd}
\end{table}

%==============================================================
\section{Nacimiento de SQL: De SEQUEL a SQL}
%==============================================================

A raíz de la publicación de Codd, IBM inició el proyecto
\textit{System R} en 1973 en sus laboratorios de San José,
California, con el objetivo de demostrar la viabilidad práctica del
modelo relacional~\cite{chamberlin1974}.
Donald D.\ Chamberlin y Raymond F.\ Boyce diseñaron un lenguaje de
consulta denominado \textbf{SEQUEL}
(\textit{Structured English Query Language}).

El nombre fue posteriormente acortado a \textbf{SQL}
(\textit{Structured Query Language}) por razones legales de marca
registrada.
En paralelo, Michael Stonebraker y Eugene Wong desarrollaron
\textbf{INGRES} en la Universidad de California, Berkeley (1973–1977),
que introdujo su propio lenguaje de consulta llamado QUEL.

\begin{nota}
Aunque coloquialmente se pronuncia tanto ``sequel'' como
``ese-cu-ele'', el estándar ISO/IEC lo denomina formalmente
``SQL''. En este curso usaremos la pronunciación ``ese-cu-ele''
por su precisión técnica.
\end{nota}

\subsection{El primer sistema comercial: Oracle}

En 1979, la empresa Relational Software Inc.\ (hoy Oracle
Corporation), fundada por Larry Ellison, liberó el primer
sistema de gestión de bases de datos relacionales comercial:
\textbf{Oracle V2}.
Paradójicamente, Oracle llegó al mercado antes que el propio
IBM con su DB2 (1983), a pesar de que la tecnología subyacente
provenía de los laboratorios de IBM.


%==============================================================
\section{¿Qué es una Base de Datos?}
%==============================================================

\begin{definicion}{Base de Datos}
Una \textbf{base de datos} es una colección organizada, persistente
y compartible de datos relacionados entre sí, diseñada para
minimizar la redundancia y maximizar la integridad, de modo que
múltiples aplicaciones y usuarios puedan acceder a ella de manera
concurrente y controlada~\cite{date2003introduction}.
\end{definicion}

Esta definición merece descomponerse en sus atributos esenciales:

\begin{itemize}
  \item \textbf{Organizada:} Los datos siguen una estructura formal
    (esquema) que define qué tipos de información se almacenan y
    cómo se relacionan entre sí.

  \item \textbf{Persistente:} Los datos sobreviven a la terminación
    de los procesos que los crean. A diferencia de las variables en
    memoria, los datos en una base de datos permanecen aunque el
    sistema se apague.

  \item \textbf{Compartible:} Múltiples usuarios o aplicaciones
    pueden acceder a los mismos datos simultáneamente, de forma
    controlada y sin interferencias.

  \item \textbf{Mínima redundancia:} El mismo dato no se repite
    innecesariamente en múltiples lugares, lo que reduce el riesgo
    de inconsistencias (un dato cambia en un lugar pero no en otro).
\end{itemize}

\subsection{Analogía con el mundo real}

Imaginemos el sistema de una biblioteca universitaria:

\begin{ejemplo}
Una biblioteca almacena información sobre \textbf{libros},
\textbf{estudiantes} y \textbf{préstamos}.
Sin una base de datos, cada ficha de préstamo copiaría el título
completo del libro, el autor, el nombre del estudiante, su carrera,
etc.
Si un estudiante cambia su dirección, habría que actualizar decenas
de fichas físicas.

Con una base de datos relacional, el estudiante aparece
\emph{una sola vez} en la tabla \texttt{Estudiantes}.
Todos los préstamos simplemente hacen referencia a su identificador.
Un cambio de dirección se realiza en un único lugar y se refleja
automáticamente en todo el sistema.
\end{ejemplo}





\section{Consulta de Datos: \texttt{SELECT}}

La instrucción \texttt{SELECT} es, sin duda, la más utilizada y la
más rica en posibilidades de todo el lenguaje SQL.
Su propósito es \textbf{consultar} datos: extraer filas de una o
varias tablas, aplicarles filtros, transformaciones y ordenamientos,
y devolver un resultado tabular.
A diferencia de \texttt{INSERT}, \texttt{UPDATE} o \texttt{DELETE},
\texttt{SELECT} es una operación de \textbf{sólo lectura}: nunca
modifica los datos almacenados.

%--------------------------------------------------------------
\subsection{Sintaxis completa y cláusulas}
%--------------------------------------------------------------

La sintaxis formal de \texttt{SELECT} en MySQL/SQL estándar es la
siguiente.
Las partes entre corchetes \texttt{[ ]} son opcionales; las palabras
en mayúsculas son palabras reservadas de SQL:

\begin{lstlisting}[caption={Sintaxis completa de SELECT.}]
SELECT  [ALL | DISTINCT | DISTINCTROW]
        lista_de_columnas | *
FROM    nombre_tabla  [[AS] alias_tabla]
[WHERE      condicion]
[GROUP BY   columna [ASC | DESC], ...]
[HAVING     condicion_de_grupo]
[ORDER BY   columna [ASC | DESC], ...]
[LIMIT      numero_de_filas  [OFFSET desplazamiento]];
\end{lstlisting}

Cada cláusula tiene un rol preciso dentro de la consulta.
La \cref{tab:clausulas-select} resume su función y el orden en que
el SGBD las \textit{evalúa} internamente (que no siempre coincide
con el orden en que se escriben):

\begin{table}[H]
\centering
\small
\renewcommand{\arraystretch}{1.4}
\begin{tabularx}{\linewidth}{@{} l c X @{}}
\toprule
\textbf{Cláusula} & \textbf{Orden eval.} & \textbf{Función} \\
\midrule
\texttt{FROM}     & 1 & Determina la(s) tabla(s) fuente de datos. \\
\texttt{WHERE}    & 2 & Filtra filas \textit{antes} de agrupar. \\
\texttt{GROUP BY} & 3 & Agrupa filas con valores iguales en las columnas indicadas. \\
\texttt{HAVING}   & 4 & Filtra grupos \textit{después} de agrupar. \\
\texttt{SELECT}   & 5 & Proyecta (selecciona) las columnas del resultado. \\
\texttt{DISTINCT} & 6 & Elimina filas duplicadas del resultado proyectado. \\
\texttt{ORDER BY} & 7 & Ordena el resultado final. \\
\texttt{LIMIT}    & 8 & Recorta el número de filas devueltas. \\
\bottomrule
\end{tabularx}
\caption{Cláusulas de \texttt{SELECT} y su orden de evaluación interno.}
\label{tab:clausulas-select}
\end{table}

\begin{nota}
El orden de \textbf{escritura} de las cláusulas es fijo y debe
respetarse exactamente como aparece en la sintaxis de arriba.
Sin embargo, el SGBD las \textbf{evalúa} en un orden diferente
(véase \cref{tab:clausulas-select}).
Esto explica, por ejemplo, por qué no se puede usar un alias
definido en \texttt{SELECT} dentro de la cláusula \texttt{WHERE}:
\texttt{WHERE} se evalúa \textit{antes} de que exista el alias.
\end{nota}

%--------------------------------------------------------------
\subsection{Tabla de ejemplo}
%--------------------------------------------------------------

Para ilustrar todos los ejemplos de esta sección usaremos una tabla
llamada \texttt{empleados} con la siguiente estructura y datos:

\begin{lstlisting}[caption={Creación e inicialización de la tabla \texttt{empleados}.}]
CREATE TABLE empleados (
    id          INT          PRIMARY KEY AUTO_INCREMENT,
    nombre      VARCHAR(60)  NOT NULL,
    departamento VARCHAR(40) NOT NULL,
    cargo       VARCHAR(40)  NOT NULL,
    salario     DECIMAL(10,2) NOT NULL,
    fecha_ingreso DATE        NOT NULL
);

INSERT INTO empleados (nombre, departamento, cargo, salario, fecha_ingreso) VALUES
('Ana García',      'Ventas',     'Gerente',    55000.00, '2018-03-15'),
('Luis Herrera',    'Ventas',     'Analista',   32000.00, '2020-07-01'),
('Sofía Ramírez',   'TI',         'Desarrolladora', 48000.00, '2019-11-20'),
('Carlos Mendoza',  'TI',         'DBA',        51000.00, '2017-05-10'),
('Elena Torres',    'RRHH',       'Coordinadora', 38000.00, '2021-02-28'),
('Marco Villanueva','Ventas',     'Analista',   31000.00, '2022-09-05'),
('Patricia Luna',   'TI',         'Desarrolladora', 47000.00, '2020-04-14'),
('Roberto Salas',   'RRHH',       'Gerente',    52000.00, '2016-08-22'),
('Diana Castillo',  'Finanzas',   'Analista',   39000.00, '2023-01-10'),
('Javier Moreno',   'Finanzas',   'Gerente',    58000.00, '2015-06-30');
\end{lstlisting}

%--------------------------------------------------------------
\subsection{Proyección de columnas: \texttt{SELECT *} y columnas específicas}
%--------------------------------------------------------------
\clearpage
La \textbf{proyección} es la operación de elegir qué columnas
aparecerán en el resultado.
El asterisco \texttt{*} es un comodín que significa
\textit{``todas las columnas''}.

\begin{lstlisting}[caption={Seleccionar todas las columnas.}]
-- Devuelve todas las columnas de todos los empleados
SELECT *
FROM empleados;
\end{lstlisting}

\begin{lstlisting}[caption={Seleccionar columnas específicas.}]
-- Sólo nombre, departamento y salario
SELECT nombre, departamento, salario
FROM empleados;
\end{lstlisting}

\begin{nota}
En producción se recomienda \textbf{evitar} \texttt{SELECT *} porque:
(1) transfiere columnas innecesarias por la red,
(2) si alguien agrega una columna sensible a la tabla, la consulta
la expondrá automáticamente,
(3) dificulta la legibilidad del código.
Úsalo sólo para exploración interactiva.
\end{nota}

También es posible incluir \textbf{expresiones} en la lista de
columnas: literales, operaciones aritméticas o llamadas a funciones:

\begin{lstlisting}[caption={Columnas calculadas y literales en SELECT.}]
SELECT
    nombre,
    salario,
    salario * 12          AS salario_anual,
    salario * 0.10        AS bono_estimado,
    'Activo'              AS estado   -- literal de texto
FROM empleados;
\end{lstlisting}

%--------------------------------------------------------------
\subsection{Operadores aritméticos}
%--------------------------------------------------------------

Dentro de \texttt{SELECT} (y también en \texttt{WHERE}) se pueden
usar operadores aritméticos estándar sobre columnas numéricas:

\begin{table}[H]
\centering
\small
\renewcommand{\arraystretch}{1.4}
\begin{tabularx}{\linewidth}{@{} c l X @{}}
\toprule
\textbf{Op.} & \textbf{Nombre} & \textbf{Ejemplo} \\
\midrule
\texttt{+}  & Suma            & \texttt{salario + 5000} \\
\texttt{-}  & Resta           & \texttt{salario - descuento} \\
\texttt{*}  & Multiplicación  & \texttt{salario * 1.15} \\
\texttt{/}  & División        & \texttt{salario / 14} \\
\texttt{\%} & Módulo (resto)  & \texttt{id \% 2} (par/impar) \\
\bottomrule
\end{tabularx}
\caption{Operadores aritméticos en SQL.}
\label{tab:operadores-aritmeticos}
\end{table}

\begin{lstlisting}[caption={Uso de operadores aritméticos.}]
SELECT
    nombre,
    salario,
    salario * 1.15        AS salario_con_aumento,
    (salario / 30)        AS salario_diario,
    salario - 5000        AS salario_base
FROM empleados
WHERE departamento = 'TI';
\end{lstlisting}

%--------------------------------------------------------------
\subsection{Operadores de comparación}
%--------------------------------------------------------------

Los operadores de comparación se usan principalmente en la cláusula
\texttt{WHERE} para filtrar filas.
Devuelven un valor booleano (\texttt{TRUE} o \texttt{FALSE}):

\begin{table}[H]
\centering
\small
\renewcommand{\arraystretch}{1.4}
\begin{tabularx}{\linewidth}{@{} c l X @{}}
\toprule
\textbf{Operador} & \textbf{Significado} & \textbf{Ejemplo} \\
\midrule
\texttt{=}    & Igual a              & \texttt{departamento = 'Ventas'} \\
\texttt{<>} o \texttt{!=} & Distinto de & \texttt{cargo <> 'Gerente'} \\
\texttt{>}    & Mayor que            & \texttt{salario > 40000} \\
\texttt{<}    & Menor que            & \texttt{salario < 35000} \\
\texttt{>=}   & Mayor o igual que    & \texttt{salario >= 50000} \\
\texttt{<=}   & Menor o igual que    & \texttt{salario <= 48000} \\
\bottomrule
\end{tabularx}
\caption{Operadores de comparación en SQL.}
\label{tab:operadores-comparacion}
\end{table}

\begin{lstlisting}[caption={Filtros con operadores de comparación.}]
-- Empleados con salario mayor a 45,000
SELECT nombre, cargo, salario
FROM   empleados
WHERE  salario > 45000;

-- Empleados que NO son del departamento de TI
SELECT nombre, departamento
FROM   empleados
WHERE  departamento <> 'TI';

-- Empleados con salario entre 30,000 y 40,000 (inclusive)
SELECT nombre, salario
FROM   empleados
WHERE  salario >= 30000
  AND  salario <= 40000;
\end{lstlisting}

\begin{ejemplo}
Una analogía útil: la cláusula \texttt{WHERE} funciona exactamente
como el filtro de una hoja de cálculo.
Si tienes una columna \texttt{salario} y aplicas el filtro
``mayor que 45,000'', sólo se muestran las filas que cumplen esa
condición. SQL hace exactamente lo mismo, pero expresado como texto.
\end{ejemplo}

%--------------------------------------------------------------
\subsection{La cláusula \texttt{WHERE}: filtrado de filas}
%--------------------------------------------------------------

\texttt{WHERE} permite especificar una \textbf{condición de filtro}
que cada fila debe satisfacer para aparecer en el resultado.
Sólo las filas para las que la condición evalúa a \texttt{TRUE}
son incluidas; las que evalúan a \texttt{FALSE} o \texttt{NULL}
son descartadas.

\begin{lstlisting}[caption={Uso básico de WHERE.}]
-- Todos los empleados del departamento de Finanzas
SELECT *
FROM   empleados
WHERE  departamento = 'Finanzas';

-- Gerentes con salario mayor o igual a 50,000
SELECT nombre, cargo, salario
FROM   empleados
WHERE  cargo    = 'Gerente'
  AND  salario >= 50000;
\end{lstlisting}

%--------------------------------------------------------------
\subsection{La cláusula \texttt{ORDER BY}: ordenamiento}
%--------------------------------------------------------------

\texttt{ORDER BY} ordena el resultado final por una o más columnas.
El orden por defecto es \textbf{ascendente} (\texttt{ASC});
se puede invertir con \textbf{descendente} (\texttt{DESC}).

\begin{lstlisting}[caption={Ordenamiento simple y compuesto.}]
-- Ordenar por salario de mayor a menor
SELECT nombre, departamento, salario
FROM   empleados
ORDER BY salario DESC;

-- Ordenar primero por departamento (A→Z)
-- y dentro de cada departamento por salario (mayor→menor)
SELECT nombre, departamento, salario
FROM   empleados
ORDER BY departamento ASC,
         salario DESC;
\end{lstlisting}

\begin{nota}
Sin \texttt{ORDER BY}, el SGBD \textbf{no garantiza} ningún orden
particular en los resultados. El orden puede cambiar entre
ejecuciones, entre versiones del SGBD o según los índices existentes.
Si el orden importa, siempre especifícalo explícitamente.
\end{nota}

%--------------------------------------------------------------
\subsection{La cláusula \texttt{LIMIT} y \texttt{OFFSET}: paginación}
%--------------------------------------------------------------

\texttt{LIMIT} restringe el número de filas devueltas.
\texttt{OFFSET} indica cuántas filas saltarse antes de comenzar
a devolver. Juntos permiten implementar \textbf{paginación}.

\begin{lstlisting}[caption={LIMIT y OFFSET para paginación.}]
-- Los 3 empleados con mayor salario
SELECT nombre, salario
FROM   empleados
ORDER BY salario DESC
LIMIT 3;

-- Página 2 (filas 4 a 6) de una lista de 3 por página
SELECT nombre, salario
FROM   empleados
ORDER BY salario DESC
LIMIT 3 OFFSET 3;

-- Sintaxis alternativa equivalente (MySQL): LIMIT offset, count
SELECT nombre, salario
FROM   empleados
ORDER BY salario DESC
LIMIT 3, 3;   -- salta 3, devuelve 3
\end{lstlisting}

%--------------------------------------------------------------
\subsection{Expresiones con \texttt{SELECT}: funciones de cadena y fecha}
%--------------------------------------------------------------

\texttt{SELECT} puede invocar funciones integradas del SGBD sobre
las columnas.
Las más utilizadas en MySQL son:

\begin{table}[H]
\centering
\small
\renewcommand{\arraystretch}{1.4}
\begin{tabularx}{\linewidth}{@{} l X @{}}
\toprule
\textbf{Función} & \textbf{Descripción} \\
\midrule
\texttt{UPPER(s)} / \texttt{LOWER(s)}  & Convierte a mayúsculas / minúsculas. \\
\texttt{LENGTH(s)}                     & Longitud en bytes de la cadena. \\
\texttt{SUBSTRING(s, inicio, n)}       & Extrae \textit{n} caracteres desde \textit{inicio}. \\
\texttt{TRIM(s)}                       & Elimina espacios al inicio y al final. \\
\texttt{CONCAT(s1, s2, ...)}           & Concatena cadenas. \\
\texttt{NOW()}                         & Fecha y hora actuales. \\
\texttt{YEAR(fecha)}                   & Extrae el año de una fecha. \\
\texttt{DATEDIFF(f1, f2)}             & Días entre dos fechas. \\
\texttt{ROUND(n, decimales)}           & Redondea un número. \\
\bottomrule
\end{tabularx}
\caption{Funciones integradas frecuentes en \texttt{SELECT}.}
\label{tab:funciones-select}
\end{table}

\begin{lstlisting}[caption={Uso de funciones de cadena, fecha y número.}]
SELECT
    UPPER(nombre)                        AS nombre_mayusculas,
    LENGTH(nombre)                       AS longitud_nombre,
    YEAR(fecha_ingreso)                  AS anio_ingreso,
    DATEDIFF(NOW(), fecha_ingreso) / 365 AS anios_en_empresa,
    ROUND(salario / 12, 2)               AS salario_mensual
FROM empleados;
\end{lstlisting}

%--------------------------------------------------------------
\subsection{Subconsultas en \texttt{SELECT}}
%--------------------------------------------------------------

Una \textbf{subconsulta} (o \textit{subquery}) es un \texttt{SELECT}
anidado dentro de otro.
En la lista de columnas se puede incluir una subconsulta
\textit{escalar} (que devuelve exactamente un valor):

\begin{lstlisting}[caption={Subconsulta escalar en la lista de columnas.}]
-- Para cada empleado, muestra cuánto gana respecto
-- al salario promedio de toda la empresa
SELECT
    nombre,
    salario,
    (SELECT AVG(salario) FROM empleados) AS promedio_empresa,
    salario - (SELECT AVG(salario) FROM empleados) AS diferencia
FROM empleados
ORDER BY diferencia DESC;
\end{lstlisting}

%--------------------------------------------------------------
\subsection{Resumen visual de la estructura de \texttt{SELECT}}
%--------------------------------------------------------------

A modo de síntesis, la siguiente figura representa gráficamente
cómo fluyen los datos a través de cada cláusula hasta producir el
resultado final:

\begin{ejemplo}
\textbf{Flujo de evaluación de una consulta SELECT}

\begin{enumerate}
  \item \texttt{FROM} — Se carga(n) la(s) tabla(s) fuente en memoria
        de trabajo.
  \item \texttt{WHERE} — Se descartan las filas que no cumplen la
        condición.
  \item \texttt{GROUP BY} — Las filas restantes se agrupan por los
        valores de las columnas indicadas.
  \item \texttt{HAVING} — Se descartan los grupos que no cumplen la
        condición de grupo.
  \item \texttt{SELECT} — De cada grupo (o fila si no hay \texttt{GROUP BY})
        se calculan y proyectan las columnas indicadas.
  \item \texttt{DISTINCT} — Se eliminan filas duplicadas del resultado
        proyectado.
  \item \texttt{ORDER BY} — El resultado se ordena.
  \item \texttt{LIMIT / OFFSET} — Se recorta la cantidad de filas
        devueltas.
\end{enumerate}
\end{ejemplo}

\begin{lstlisting}[caption={Consulta que usa todas las cláusulas.}]
-- "Los 3 departamentos donde el salario promedio
--  supera los 40,000, ordenados de mayor a menor promedio,
--  excluyendo el departamento de RRHH."
SELECT
    departamento,
    COUNT(*)          AS num_empleados,
    AVG(salario)      AS salario_promedio,
    MAX(salario)      AS salario_maximo
FROM    empleados
WHERE   departamento <> 'RRHH'          -- filtra filas
GROUP BY departamento                    -- agrupa
HAVING  AVG(salario) > 40000            -- filtra grupos
ORDER BY salario_promedio DESC           -- ordena
LIMIT   3;                               -- recorta
\end{lstlisting}
% ─── MÓDULO 4: Modificadores Parte 1 ─────────────────────
\newpage
% ─── MÓDULO 12: Consulta de datos relacionados ───────────
\newpage
\section{Consulta de Datos Relacionados}

Las consultas relacionales permiten combinar información distribuida en múltiples tablas mediante
operaciones de \textit{join} y unión. SQL ofrece distintos tipos de combinaciones según el
comportamiento deseado ante filas sin coincidencia.

\begin{nota}
  Para los ejemplos de esta sección se utilizan las tablas \texttt{clientes} y \texttt{pedidos},
  relacionadas por la columna \texttt{cliente\_id}.
\end{nota}

% ── Tablas de ejemplo ─────────────────────────────────────
\begin{table}[H]
  \centering
  \caption{Tabla \texttt{clientes}}
  \begin{tabular}{>{\ttfamily}c >{\ttfamily}l}
    \toprule
    \textbf{cliente\_id} & \textbf{nombre} \\
    \midrule
    1  & Ana García   \\
    2  & Luis Pérez   \\
    3  & María López  \\
    4  & Carlos Ruiz  \\
    \bottomrule
  \end{tabular}
\end{table}

\begin{table}[H]
  \centering
  \caption{Tabla \texttt{pedidos}}
  \begin{tabular}{>{\ttfamily}c >{\ttfamily}c >{\ttfamily}r}
    \toprule
    \textbf{pedido\_id} & \textbf{cliente\_id} & \textbf{monto} \\
    \midrule
    101 & 1 & 350.00 \\
    102 & 2 & 120.50 \\
    103 & 2 &  89.00 \\
    104 & 5 & 200.00 \\
    \bottomrule
  \end{tabular}
\end{table}

% ─────────────────────────────────────────────────────────
\subsection{\texttt{INNER JOIN}}

\begin{definicion}{INNER JOIN}
  Devuelve únicamente las filas que tienen \textbf{coincidencia en ambas tablas}.
  Las filas sin par en la tabla opuesta quedan excluidas del resultado.
\end{definicion}

\begin{center}
  \begin{tikzpicture}
    \fill[NavyBlue!30] (0,0) circle (1.2);
    \fill[ForestGreen!30] (1.6,0) circle (1.2);
    \begin{scope}
      \clip (0,0) circle (1.2);
      \fill[NavyBlue!60] (1.6,0) circle (1.2);
    \end{scope}
    \draw (0,0) circle (1.2) node[left=1.1cm] {\small\texttt{clientes}};
    \draw (1.6,0) circle (1.2) node[right=1.1cm] {\small\texttt{pedidos}};
    \node at (0.8,0) {\tiny JOIN};
  \end{tikzpicture}
\end{center}

\textbf{Sintaxis general:}

\begin{lstlisting}[style=sql, caption={Sintaxis de INNER JOIN}]
SELECT columnas
FROM tabla_izquierda
INNER JOIN tabla_derecha
    ON tabla_izquierda.clave = tabla_derecha.clave;
\end{lstlisting}

\begin{ejemplo}
  Obtener los pedidos junto con el nombre del cliente que los realizó:

  \begin{lstlisting}[style=sql, caption={INNER JOIN entre clientes y pedidos}]
SELECT c.nombre,
       p.pedido_id,
       p.monto
FROM clientes AS c
INNER JOIN pedidos AS p
    ON c.cliente_id = p.cliente_id;
  \end{lstlisting}

  \textbf{Resultado:}

  \begin{center}
    \begin{tabular}{>{\ttfamily}l >{\ttfamily}c >{\ttfamily}r}
      \toprule
      \textbf{nombre}  & \textbf{pedido\_id} & \textbf{monto} \\
      \midrule
      Ana García  & 101 & 350.00 \\
      Luis Pérez  & 102 & 120.50 \\
      Luis Pérez  & 103 &  89.00 \\
      \bottomrule
    \end{tabular}
  \end{center}

  Los clientes \texttt{María López} y \texttt{Carlos Ruiz} no aparecen porque no tienen
  pedidos. El pedido \texttt{104} tampoco aparece porque su \texttt{cliente\_id = 5}
  no existe en \texttt{clientes}.
\end{ejemplo}

\begin{nota}
  La palabra clave \texttt{INNER} es opcional; \texttt{JOIN} a secas equivale a
  \texttt{INNER JOIN} en todos los motores SQL principales.
\end{nota}

% ─────────────────────────────────────────────────────────
\subsection{\texttt{LEFT JOIN}}

\begin{definicion}{LEFT JOIN}
  Devuelve \textbf{todas las filas de la tabla izquierda} más las coincidencias de la tabla
  derecha. Cuando no existe par, las columnas de la tabla derecha se rellenan con
  \texttt{NULL}.
\end{definicion}

\begin{center}
  \begin{tikzpicture}
    \fill[NavyBlue!50] (0,0) circle (1.2);
    \fill[ForestGreen!30] (1.6,0) circle (1.2);
    \begin{scope}
      \clip (0,0) circle (1.2);
      \fill[NavyBlue!50] (1.6,0) circle (1.2);
    \end{scope}
    \draw (0,0) circle (1.2) node[left=1.1cm] {\small\texttt{clientes}};
    \draw (1.6,0) circle (1.2) node[right=1.1cm] {\small\texttt{pedidos}};
  \end{tikzpicture}
\end{center}

\textbf{Sintaxis general:}

\begin{lstlisting}[style=sql, caption={Sintaxis de LEFT JOIN}]
SELECT columnas
FROM tabla_izquierda
LEFT JOIN tabla_derecha
    ON tabla_izquierda.clave = tabla_derecha.clave;
\end{lstlisting}

\begin{ejemplo}
  Listar todos los clientes y, si tienen pedidos, mostrar su monto:

  \begin{lstlisting}[style=sql, caption={LEFT JOIN: todos los clientes con sus pedidos}]
SELECT c.nombre,
       p.pedido_id,
       p.monto
FROM clientes AS c
LEFT JOIN pedidos AS p
    ON c.cliente_id = p.cliente_id;
  \end{lstlisting}

  \textbf{Resultado:}

  \begin{center}
    \begin{tabular}{>{\ttfamily}l >{\ttfamily}c >{\ttfamily}r}
      \toprule
      \textbf{nombre}  & \textbf{pedido\_id} & \textbf{monto} \\
      \midrule
      Ana García   & 101  & 350.00 \\
      Luis Pérez   & 102  & 120.50 \\
      Luis Pérez   & 103  &  89.00 \\
      María López  & NULL &   NULL \\
      Carlos Ruiz  & NULL &   NULL \\
      \bottomrule
    \end{tabular}
  \end{center}

  Los clientes sin pedidos aparecen con \texttt{NULL} en las columnas de \texttt{pedidos}.
\end{ejemplo}

\textbf{Caso de uso habitual — detectar huérfanos:}

\begin{lstlisting}[style=sql, caption={Clientes que nunca han realizado un pedido}]
SELECT c.nombre
FROM clientes AS c
LEFT JOIN pedidos AS p
    ON c.cliente_id = p.cliente_id
WHERE p.pedido_id IS NULL;
\end{lstlisting}

% ─────────────────────────────────────────────────────────
\subsection{\texttt{RIGHT JOIN}}

\begin{definicion}{RIGHT JOIN}
  Devuelve \textbf{todas las filas de la tabla derecha} más las coincidencias de la tabla
  izquierda. Las columnas de la tabla izquierda se rellenan con \texttt{NULL} cuando no
  existe par. Es el simétrico de \texttt{LEFT JOIN}.
\end{definicion}

\begin{center}
  \begin{tikzpicture}
    \fill[NavyBlue!30] (0,0) circle (1.2);
    \fill[ForestGreen!50] (1.6,0) circle (1.2);
    \begin{scope}
      \clip (1.6,0) circle (1.2);
      \fill[ForestGreen!50] (0,0) circle (1.2);
    \end{scope}
    \draw (0,0) circle (1.2) node[left=1.1cm] {\small\texttt{clientes}};
    \draw (1.6,0) circle (1.2) node[right=1.1cm] {\small\texttt{pedidos}};
  \end{tikzpicture}
\end{center}

\textbf{Sintaxis general:}

\begin{lstlisting}[style=sql, caption={Sintaxis de RIGHT JOIN}]
SELECT columnas
FROM tabla_izquierda
RIGHT JOIN tabla_derecha
    ON tabla_izquierda.clave = tabla_derecha.clave;
\end{lstlisting}

\begin{ejemplo}
  Mostrar todos los pedidos e incluir el nombre del cliente si existe:

  \begin{lstlisting}[style=sql, caption={RIGHT JOIN: todos los pedidos con datos del cliente}]
SELECT c.nombre,
       p.pedido_id,
       p.monto
FROM clientes AS c
RIGHT JOIN pedidos AS p
    ON c.cliente_id = p.cliente_id;
  \end{lstlisting}

  \textbf{Resultado:}

  \begin{center}
    \begin{tabular}{>{\ttfamily}l >{\ttfamily}c >{\ttfamily}r}
      \toprule
      \textbf{nombre}  & \textbf{pedido\_id} & \textbf{monto} \\
      \midrule
      Ana García  & 101 & 350.00 \\
      Luis Pérez  & 102 & 120.50 \\
      Luis Pérez  & 103 &  89.00 \\
      NULL        & 104 & 200.00 \\
      \bottomrule
    \end{tabular}
  \end{center}

  El pedido \texttt{104} aparece con \texttt{NULL} en \texttt{nombre} porque su
  \texttt{cliente\_id = 5} no existe en la tabla \texttt{clientes}.
\end{ejemplo}

\begin{nota}
  En la práctica, \texttt{RIGHT JOIN} se usa con menos frecuencia que \texttt{LEFT JOIN}.
  Cualquier consulta con \texttt{RIGHT JOIN} puede reescribirse como un \texttt{LEFT JOIN}
  intercambiando el orden de las tablas:
  \begin{center}
    \texttt{A RIGHT JOIN B} $\equiv$ \texttt{B LEFT JOIN A}
  \end{center}
\end{nota}

% ─────────────────────────────────────────────────────────
\subsection{\texttt{UNION}}

\begin{definicion}{UNION}
  Combina verticalmente los resultados de dos o más sentencias \texttt{SELECT}.
  \texttt{UNION} elimina filas duplicadas; \texttt{UNION ALL} las conserva.
  Ambas consultas deben tener el \textbf{mismo número de columnas} y tipos de datos
  compatibles.
\end{definicion}

\textbf{Sintaxis general:}

\begin{lstlisting}[style=sql, caption={Sintaxis de UNION / UNION ALL}]
SELECT columna1, columna2, ...
FROM tabla_a
UNION [ALL]
SELECT columna1, columna2, ...
FROM tabla_b;
\end{lstlisting}

\begin{ejemplo}
  Obtener una lista unificada de correos de dos tablas distintas (\texttt{clientes} y
  \texttt{proveedores}), sin duplicados:

  \begin{lstlisting}[style=sql, caption={UNION de correos de clientes y proveedores}]
SELECT nombre, correo, 'cliente'   AS origen
FROM clientes
UNION
SELECT nombre, correo, 'proveedor' AS origen
FROM proveedores
ORDER BY nombre;
  \end{lstlisting}
\end{ejemplo}

\begin{ejemplo}
  Comparación entre \texttt{UNION} y \texttt{UNION ALL}:

  \begin{lstlisting}[style=sql, caption={Diferencia entre UNION y UNION ALL}]
-- Elimina duplicados (más lento: realiza ORDER + DISTINCT internamente)
SELECT ciudad FROM clientes
UNION
SELECT ciudad FROM proveedores;

-- Conserva duplicados (más rápido: no realiza deduplicación)
SELECT ciudad FROM clientes
UNION ALL
SELECT ciudad FROM proveedores;
  \end{lstlisting}
\end{ejemplo}

\begin{nota}
  Los nombres de las columnas del resultado se toman del \textbf{primer} \texttt{SELECT}.
  La cláusula \texttt{ORDER BY} solo puede aparecer \textbf{una vez}, al final de toda la
  sentencia, y afecta al conjunto combinado completo.
\end{nota}

% ── Tabla resumen ─────────────────────────────────────────
\subsection*{Resumen comparativo}

\begin{table}[H]
  \centering
  \caption{Comparación de operaciones de combinación en SQL}
  \renewcommand{\arraystretch}{1.3}
  \rowcolors{2}{gray!10}{white}
  \begin{tabularx}{\textwidth}{>{\ttfamily}l X X}
    \toprule
    \textbf{Operación}  & \textbf{Filas incluidas}                                   & \textbf{Uso típico}                         \\
    \midrule
    INNER JOIN  & Solo las que tienen coincidencia en ambas tablas             & Datos relacionados que siempre existen      \\
    LEFT JOIN   & Todas las de la tabla izquierda + coincidencias              & Encontrar registros sin par (huérfanos)     \\
    RIGHT JOIN  & Todas las de la tabla derecha + coincidencias                & Simétrico de LEFT JOIN                      \\
    UNION       & Filas de ambos SELECT, sin duplicados                        & Consolidar datos de tablas similares        \\
    UNION ALL   & Filas de ambos SELECT, con duplicados                        & Igual que UNION pero más eficiente          \\
    \bottomrule
  \end{tabularx}
\end{table}


\section{Modificadores: Parte 1}

Los \textbf{modificadores} son cláusulas y operadores que se
combinan con \texttt{SELECT} para refinar, filtrar y ordenar los
resultados de una consulta.
En esta sección cubrimos los seis modificadores fundamentales que
todo programador SQL debe dominar desde el primer día.
Todos los ejemplos siguen usando la tabla \texttt{empleados}
introducida en la sección anterior.

%--------------------------------------------------------------
\subsection{\texttt{DISTINCT}}
%--------------------------------------------------------------

Por defecto, \texttt{SELECT} devuelve \textbf{todas} las filas que
satisfacen la consulta, incluyendo duplicados.
La palabra clave \texttt{DISTINCT} instruye al SGBD a eliminar las
filas repetidas del resultado, conservando sólo valores únicos.

\begin{definicion}{\texttt{DISTINCT}}
Modificador que se coloca inmediatamente después de \texttt{SELECT}
y elimina del resultado todas las filas cuyo conjunto de valores en
las columnas proyectadas sea idéntico a otra fila ya presente.
\end{definicion}

\begin{lstlisting}[caption={DISTINCT sobre una columna.}]
-- ¿Qué departamentos existen en la empresa?
-- Sin DISTINCT aparecerían 10 filas (una por empleado)
SELECT DISTINCT departamento
FROM   empleados;

-- Resultado:
-- Ventas
-- TI
-- RRHH
-- Finanzas
\end{lstlisting}

\begin{lstlisting}[caption={DISTINCT sobre múltiples columnas.}]
-- Combinaciones únicas de departamento + cargo
SELECT DISTINCT departamento, cargo
FROM   empleados
ORDER BY departamento, cargo;
\end{lstlisting}

\begin{nota}
\texttt{DISTINCT} opera sobre la \textit{combinación completa} de
todas las columnas proyectadas, no columna por columna.
En el ejemplo anterior, \texttt{('TI', 'Desarrolladora')} es una
combinación única aunque haya dos desarrolladoras en TI, porque
ambas tienen exactamente los mismos valores en esas dos columnas
— en ese caso sí se eliminará el duplicado.
\end{nota}

%--------------------------------------------------------------
\subsection{\texttt{WHERE}}
%--------------------------------------------------------------

La cláusula \texttt{WHERE} es el mecanismo de \textbf{filtrado por
fila} de SQL.
Evalúa una condición booleana para cada fila de la tabla fuente;
sólo las filas donde la condición resulta \texttt{TRUE} pasan al
resultado.

\begin{definicion}{\texttt{WHERE}}
Cláusula que restringe las filas devueltas por una consulta a
aquellas que satisfacen una expresión booleana dada.
Se evalúa \textit{antes} de \texttt{GROUP BY}, \texttt{HAVING} y
\texttt{SELECT}.
\end{definicion}

\begin{lstlisting}[caption={Filtros básicos con WHERE.}]
-- Empleados del departamento de TI
SELECT nombre, cargo, salario
FROM   empleados
WHERE  departamento = 'TI';

-- Empleados con salario mayor a 45,000
SELECT nombre, salario
FROM   empleados
WHERE  salario > 45000;

-- Empleados que ingresaron después del 1 de enero de 2020
SELECT nombre, fecha_ingreso
FROM   empleados
WHERE  fecha_ingreso > '2020-01-01';
\end{lstlisting}

\begin{ejemplo}
Imagina la tabla como un tamiz: cada fila cae por los agujeros a
menos que cumpla la condición de \texttt{WHERE}.
Las que no la cumplen simplemente desaparecen del resultado.
El dato original en la tabla \textbf{no se modifica}; sólo se
filtra lo que se muestra.
\end{ejemplo}

%--------------------------------------------------------------
\subsection{\texttt{ORDER BY}}
%--------------------------------------------------------------

\texttt{ORDER BY} define el \textbf{criterio de ordenamiento} del
resultado final.
Sin él, el SGBD no garantiza ningún orden particular entre
ejecuciones.

\begin{lstlisting}[caption={Orden ascendente y descendente.}]
-- Orden ascendente (A → Z, menor → mayor): por defecto
SELECT nombre, salario
FROM   empleados
ORDER BY salario ASC;   -- ASC es opcional, es el valor por defecto

-- Orden descendente (mayor → menor)
SELECT nombre, salario
FROM   empleados
ORDER BY salario DESC;
\end{lstlisting}

\begin{lstlisting}[caption={Ordenamiento compuesto por múltiples columnas.}]
-- Primero por departamento (A→Z),
-- dentro de cada departamento por salario (mayor→menor)
SELECT nombre, departamento, salario
FROM   empleados
ORDER BY departamento ASC,
         salario      DESC;
\end{lstlisting}

\begin{lstlisting}[caption={ORDER BY usando posición numérica de columna.}]
-- Se puede referenciar la columna por su posición en SELECT
-- (poco recomendado: hace el código menos legible)
SELECT nombre, departamento, salario
FROM   empleados
ORDER BY 3 DESC;   -- 3 = tercera columna = salario
\end{lstlisting}

\begin{nota}
El uso de posición numérica en \texttt{ORDER BY} está permitido
pero se desaconseja en código de producción: si el orden de las
columnas en \texttt{SELECT} cambia, el ordenamiento cambia sin
ningún aviso de error.
\end{nota}

%--------------------------------------------------------------
\subsection{\texttt{LIKE}}
%--------------------------------------------------------------

\texttt{LIKE} permite comparar una columna de texto contra un
\textbf{patrón}, usando caracteres comodín.
Es la herramienta de búsqueda de texto parcial más básica de SQL.

\begin{table}[H]
\centering
\small
\renewcommand{\arraystretch}{1.4}
\begin{tabularx}{\linewidth}{@{} c l X @{}}
\toprule
\textbf{Comodín} & \textbf{Significado} & \textbf{Ejemplo} \\
\midrule
\texttt{\%} & Cero o más caracteres cualesquiera
            & \texttt{'Ana\%'} coincide con \texttt{Ana}, \texttt{Analista}, \texttt{Anabel} \\
\texttt{\_} & Exactamente un carácter cualquiera
            & \texttt{'\_na'} coincide con \texttt{Ana}, \texttt{Ina}, \texttt{Una} \\
\bottomrule
\end{tabularx}
\caption{Comodines del operador \texttt{LIKE}.}
\label{tab:comodines-like}
\end{table}

\begin{lstlisting}[caption={Patrones con LIKE.}]
-- Nombres que empiezan con 'A'
SELECT nombre FROM empleados
WHERE  nombre LIKE 'A%';

-- Nombres que terminan con 'ez'
SELECT nombre FROM empleados
WHERE  nombre LIKE '%ez';

-- Nombres que contienen 'ar' en cualquier posición
SELECT nombre FROM empleados
WHERE  nombre LIKE '%ar%';

-- Nombres cuya segunda letra es 'o'
SELECT nombre FROM empleados
WHERE  nombre LIKE '_o%';

-- Nombres que NO contienen 'a' (NOT LIKE)
SELECT nombre FROM empleados
WHERE  nombre NOT LIKE '%a%';
\end{lstlisting}

\begin{nota}
En MySQL, \texttt{LIKE} es \textbf{insensible a mayúsculas/minúsculas}
por defecto cuando la columna usa un \textit{collation} insensible
(\texttt{utf8\_general\_ci}, \texttt{utf8mb4\_unicode\_ci}, etc.).
Si necesitas distinguir mayúsculas usa \texttt{LIKE BINARY} o cambia
el collation de la columna.
\end{nota}

%--------------------------------------------------------------
\subsection{\texttt{AND}, \texttt{OR}, \texttt{NOT}}
%--------------------------------------------------------------

Los operadores lógicos permiten \textbf{combinar múltiples
condiciones} dentro de una cláusula \texttt{WHERE}.

\begin{table}[H]
\centering
\small
\renewcommand{\arraystretch}{1.5}
\begin{tabularx}{\linewidth}{@{} l X X @{}}
\toprule
\textbf{Operador} & \textbf{Resultado TRUE cuando\ldots} & \textbf{Precedencia} \\
\midrule
\texttt{NOT} & La condición que sigue es \texttt{FALSE}
             & Mayor (se evalúa primero) \\
\texttt{AND} & \textbf{Ambas} condiciones son \texttt{TRUE}
             & Media \\
\texttt{OR}  & \textbf{Al menos una} condición es \texttt{TRUE}
             & Menor (se evalúa al final) \\
\bottomrule
\end{tabularx}
\caption{Operadores lógicos y su precedencia.}
\label{tab:operadores-logicos}
\end{table}

\begin{lstlisting}[caption={Uso de AND, OR y NOT.}]
-- AND: empleados de TI con salario mayor a 47,000
SELECT nombre, departamento, salario
FROM   empleados
WHERE  departamento = 'TI'
  AND  salario > 47000;

-- OR: empleados de Ventas O de Finanzas
SELECT nombre, departamento
FROM   empleados
WHERE  departamento = 'Ventas'
  OR   departamento = 'Finanzas';

-- NOT: todos los departamentos excepto RRHH
SELECT nombre, departamento
FROM   empleados
WHERE  NOT departamento = 'RRHH';
-- Equivalente: WHERE departamento <> 'RRHH'
\end{lstlisting}

\begin{lstlisting}[caption={Combinación de operadores y uso de paréntesis.}]
-- Gerentes de TI O analistas de Ventas con salario > 30,000
-- Los paréntesis son ESENCIALES aquí
SELECT nombre, departamento, cargo, salario
FROM   empleados
WHERE  (departamento = 'TI'     AND cargo = 'Gerente')
  OR   (departamento = 'Ventas' AND cargo = 'Analista'
                                AND salario > 30000);
\end{lstlisting}

\begin{nota}
\texttt{AND} tiene mayor precedencia que \texttt{OR}, igual que la
multiplicación tiene mayor precedencia que la suma en aritmética.
Sin paréntesis, \texttt{A OR B AND C} se interpreta como
\texttt{A OR (B AND C)}, lo que puede producir resultados
inesperados.
\textbf{Usa siempre paréntesis} cuando combines \texttt{AND} y
\texttt{OR} en la misma condición.
\end{nota}

%--------------------------------------------------------------
\subsection{\texttt{LIMIT}}
%--------------------------------------------------------------

\texttt{LIMIT} restringe el número máximo de filas que la consulta
devuelve.
Combinado con \texttt{OFFSET} (o la sintaxis abreviada de MySQL)
permite implementar \textbf{paginación}.

\begin{lstlisting}[caption={LIMIT básico.}]
-- Los 5 empleados con mayor salario
SELECT nombre, salario
FROM   empleados
ORDER BY salario DESC
LIMIT 5;
\end{lstlisting}

\begin{lstlisting}[caption={LIMIT con OFFSET para paginación.}]
-- Página 1: primeras 3 filas (OFFSET 0, implícito)
SELECT nombre, departamento
FROM   empleados
ORDER BY nombre ASC
LIMIT 3;

-- Página 2: filas 4, 5 y 6
SELECT nombre, departamento
FROM   empleados
ORDER BY nombre ASC
LIMIT 3 OFFSET 3;

-- Página 3: filas 7, 8 y 9
SELECT nombre, departamento
FROM   empleados
ORDER BY nombre ASC
LIMIT 3 OFFSET 6;
\end{lstlisting}

\begin{lstlisting}[caption={Sintaxis abreviada MySQL: LIMIT offset, cantidad.}]
-- Equivalente a LIMIT 3 OFFSET 6
SELECT nombre, departamento
FROM   empleados
ORDER BY nombre ASC
LIMIT 6, 3;    -- primero el offset, luego la cantidad
\end{lstlisting}

\begin{ejemplo}
La paginación es imprescindible en aplicaciones web.
Si una tabla tiene 500,000 registros y una página muestra 20,
la consulta para la página 3 sería:
\texttt{LIMIT 20 OFFSET 40}.
Sin \texttt{LIMIT}, traer 500,000 filas al navegador colapsaría
tanto la base de datos como la red.
\end{ejemplo}

\begin{nota}
\texttt{LIMIT} sin \texttt{ORDER BY} produce resultados
\textbf{no deterministas}: en dos ejecuciones distintas podrían
devolverse filas diferentes.
Siempre acompáñalo de \texttt{ORDER BY} para obtener páginas
consistentes.
\end{nota}



% ─── MÓDULO 5: Modificadores Parte 2 ─────────────────────
\newpage

\section{Modificadores: Parte 2}

En esta sección cubrimos los modificadores avanzados de consulta:
funciones de agregación, operadores de rango y lista, alias,
concatenación, agrupamiento, lógica condicional y manejo de nulos.
Todos son componentes esenciales para construir consultas analíticas
y reportes sobre los datos.


%--------------------------------------------------------------
\subsection{Valores \texttt{NULL}}
%--------------------------------------------------------------

\texttt{NULL} no es un valor como \texttt{0} o \texttt{''};
es la \textbf{ausencia de valor}.
Representa información desconocida o no aplicable, y su
comportamiento en operaciones es especial:
cualquier operación aritmética o comparación con \texttt{NULL}
devuelve \texttt{NULL} (ni \texttt{TRUE} ni \texttt{FALSE}).

\begin{definicion}{\texttt{NULL}}
Marcador especial que indica la ausencia de un valor en una columna.
No es equivalente a cero, a una cadena vacía ni a \texttt{FALSE}.
Para comprobar si un valor es \texttt{NULL} se usan los operadores
\texttt{IS NULL} e \texttt{IS NOT NULL}; el operador \texttt{=}
nunca devuelve \texttt{TRUE} al comparar con \texttt{NULL}.
\end{definicion}

\begin{lstlisting}[caption={Comportamiento de NULL y cómo detectarlo.}]
-- INCORRECTO: esta condición nunca devuelve filas
-- porque NULL = NULL es NULL (no TRUE)
SELECT * FROM empleados WHERE telefono = NULL;   -- ¡MAL!

-- CORRECTO: usar IS NULL
SELECT * FROM empleados WHERE telefono IS NULL;


-- Empleados que SÍ tienen teléfono registrado
SELECT * FROM empleados WHERE telefono IS NOT NULL;

-- NULL en aritmética: el resultado es siempre NULL
SELECT nombre, salario + NULL AS resultado  -- → NULL
FROM   empleados;

-- NULL en ORDER BY: MySQL los coloca al INICIO en ASC
-- y al FINAL en DESC (varía entre SGBD)
SELECT nombre, telefono
FROM   empleados
ORDER BY telefono ASC;
\end{lstlisting}

\begin{ejemplo}
Supón que agregas una columna \texttt{telefono} a la tabla
\texttt{empleados} sin indicar \texttt{NOT NULL}.
Los registros existentes quedarán con \texttt{NULL} en esa columna,
no con \texttt{0} ni con \texttt{''}.
Si después calculas \texttt{COUNT(telefono)}, esos registros
\textit{no} se contarán, porque \texttt{COUNT} ignora \texttt{NULL}s.
\end{ejemplo}

%--------------------------------------------------------------
\subsection{\texttt{MIN} y \texttt{MAX}}
%--------------------------------------------------------------

\texttt{MIN} y \texttt{MAX} son \textbf{funciones de agregación}
que devuelven el valor mínimo y máximo de una columna dentro de
un conjunto de filas.
Funcionan sobre números, fechas y cadenas de texto.

\begin{lstlisting}[caption={MIN y MAX sobre números y fechas.}]
-- Salario más bajo y más alto de toda la empresa
SELECT MIN(salario) AS salario_minimo,
       MAX(salario) AS salario_maximo
FROM   empleados;

-- MIN y MAX por departamento
SELECT departamento,
       MIN(salario) AS minimo,
       MAX(salario) AS maximo,
       MAX(salario) - MIN(salario) AS rango_salarial
FROM   empleados
GROUP BY departamento
ORDER BY rango_salarial DESC;

-- Fecha de ingreso más antigua y más reciente
SELECT MIN(fecha_ingreso) AS primer_ingreso,
       MAX(fecha_ingreso) AS ultimo_ingreso
FROM   empleados;
\end{lstlisting}

\begin{nota}
\texttt{MIN} y \texttt{MAX} \textbf{ignoran} los valores
\texttt{NULL} automáticamente.
Si todas las filas de un grupo son \texttt{NULL}, el resultado
de la función es \texttt{NULL}.
\end{nota}

%--------------------------------------------------------------
\subsection{\texttt{COUNT}}
%--------------------------------------------------------------

\texttt{COUNT} cuenta filas o valores no nulos.
Tiene tres variantes con comportamientos distintos:

\begin{lstlisting}[caption={Variantes de COUNT.}]
-- COUNT(*): cuenta TODAS las filas, incluyendo NULLs
SELECT COUNT(*) AS total_empleados
FROM   empleados;
-- Resultado: 10

-- COUNT(columna): cuenta sólo las filas donde
-- esa columna NO es NULL
SELECT COUNT(telefono) AS empleados_con_telefono
FROM   empleados;
-- Si 3 tienen NULL en telefono → resultado: 7

-- COUNT(DISTINCT columna): cuenta valores únicos no nulos
SELECT COUNT(DISTINCT departamento) AS num_departamentos
FROM   empleados;
-- Resultado: 4

-- COUNT con GROUP BY: total por grupo
SELECT departamento, COUNT(*) AS num_empleados
FROM   empleados
GROUP BY departamento
ORDER BY num_empleados DESC;
\end{lstlisting}

%--------------------------------------------------------------
\subsection{\texttt{SUM}}
%--------------------------------------------------------------

\texttt{SUM} suma todos los valores no nulos de una columna
numérica dentro del conjunto de filas evaluado.

\begin{lstlisting}[caption={SUM: suma total y por grupo.}]
-- Masa salarial total de la empresa
SELECT SUM(salario) AS masa_salarial_total
FROM   empleados;

-- Masa salarial por departamento
SELECT departamento,
       SUM(salario)  AS masa_salarial,
       COUNT(*)      AS num_empleados
FROM   empleados
GROUP BY departamento
ORDER BY masa_salarial DESC;

-- SUM condicional usando CASE (suma sólo gerentes)
SELECT SUM(CASE WHEN cargo = 'Gerente'
                THEN salario ELSE 0 END) AS total_gerentes
FROM   empleados;
\end{lstlisting}

%--------------------------------------------------------------
\subsection{\texttt{AVG}}
%--------------------------------------------------------------

\texttt{AVG} calcula el \textbf{promedio aritmético} de los valores
no nulos de una columna numérica.
Equivale a \texttt{SUM(col) / COUNT(col)}, pero ignorando
\texttt{NULL}s en ambos términos.

\begin{lstlisting}[caption={AVG: promedio general y por grupo.}]
-- Salario promedio de toda la empresa
SELECT ROUND(AVG(salario), 2) AS salario_promedio
FROM   empleados;

-- Promedio por departamento, con indicador sobre/bajo media global
SELECT departamento,
       ROUND(AVG(salario), 2) AS promedio_dpto,
       ROUND(AVG(salario) - (SELECT AVG(salario)
                             FROM empleados), 2) AS diferencia_global
FROM   empleados
GROUP BY departamento
ORDER BY promedio_dpto DESC;
\end{lstlisting}

\begin{nota}
\texttt{AVG} ignora los \texttt{NULL}s: si una columna tiene 10
filas pero 3 son \texttt{NULL}, el promedio se calcula sobre los
7 valores restantes, no dividiendo entre 10.
Si quieres tratar los \texttt{NULL}s como cero, usa
\texttt{AVG(COALESCE(columna, 0))}.
\end{nota}

%--------------------------------------------------------------
\subsection{\texttt{IN}}
%--------------------------------------------------------------

\texttt{IN} verifica si el valor de una columna pertenece a una
\textbf{lista de valores} especificada.
Es equivalente a una cadena de condiciones \texttt{OR} sobre la
misma columna, pero más legible y compacto.

\begin{lstlisting}[caption={IN con lista de valores literales.}]
-- Empleados de Ventas, TI o Finanzas
-- Equivalente: WHERE departamento = 'Ventas'
--              OR departamento = 'TI'
--              OR departamento = 'Finanzas'
SELECT nombre, departamento
FROM   empleados
WHERE  departamento IN ('Ventas', 'TI', 'Finanzas');

-- NOT IN: departamentos excluidos
SELECT nombre, departamento
FROM   empleados
WHERE  departamento NOT IN ('RRHH');

-- IN con subconsulta: departamentos que tienen gerente
SELECT nombre, departamento
FROM   empleados
WHERE  departamento IN (
    SELECT DISTINCT departamento
    FROM   empleados
    WHERE  cargo = 'Gerente'
);
\end{lstlisting}

\begin{nota}
\textbf{Precaución con \texttt{NOT IN} y \texttt{NULL}s:}
si la lista de \texttt{NOT IN} contiene un \texttt{NULL},
la condición completa devuelve \texttt{NULL} (ninguna fila),
porque SQL no puede determinar si el valor es distinto de algo
desconocido.
En subconsultas, usa \texttt{NOT EXISTS} en lugar de \texttt{NOT IN}
si los resultados pueden contener \texttt{NULL}s.
\end{nota}

%--------------------------------------------------------------
\subsection{\texttt{BETWEEN}}
%--------------------------------------------------------------

\texttt{BETWEEN} filtra filas cuyos valores se encuentran dentro
de un \textbf{rango inclusivo} (incluye los extremos).
Funciona con números, fechas y cadenas.

\begin{lstlisting}[caption={BETWEEN con números y fechas.}]
-- Empleados con salario entre 35,000 y 52,000 (inclusive)
-- Equivalente: WHERE salario >= 35000 AND salario <= 52000
SELECT nombre, salario
FROM   empleados
WHERE  salario BETWEEN 35000 AND 52000;

-- NOT BETWEEN: fuera del rango
SELECT nombre, salario
FROM   empleados
WHERE  salario NOT BETWEEN 35000 AND 52000;

-- BETWEEN con fechas
SELECT nombre, fecha_ingreso
FROM   empleados
WHERE  fecha_ingreso BETWEEN '2019-01-01' AND '2021-12-31';

-- BETWEEN con cadenas (orden alfabético)
SELECT nombre
FROM   empleados
WHERE  nombre BETWEEN 'C' AND 'M';  -- nombres de C a M
\end{lstlisting}

%--------------------------------------------------------------
\subsection{\texttt{ALIAS} (\texttt{AS})}
%--------------------------------------------------------------

Un \textbf{alias} es un nombre temporal que se asigna a una columna
o tabla dentro de una consulta.
Mejora la legibilidad del resultado y es indispensable al trabajar
con expresiones calculadas.

\begin{lstlisting}[caption={Alias de columna con AS.}]
SELECT
    nombre                          AS empleado,
    salario                         AS sueldo_base,
    salario * 12                    AS ingreso_anual,
    ROUND(salario / 30, 2)          AS sueldo_diario,
    CONCAT(nombre, ' (', cargo, ')') AS descripcion
FROM empleados;
\end{lstlisting}

\begin{lstlisting}[caption={Alias de tabla: útil con JOINs y subconsultas.}]
-- Alias de tabla: simplifica nombres largos
SELECT e.nombre, e.departamento, e.salario
FROM   empleados AS e
WHERE  e.salario > 40000;

-- La palabra AS es opcional (funciona igual sin ella)
SELECT e.nombre
FROM   empleados e
WHERE  e.cargo = 'Gerente';
\end{lstlisting}

\begin{nota}
Los alias de columna definidos en \texttt{SELECT} \textbf{no pueden
usarse} en la cláusula \texttt{WHERE} (porque \texttt{WHERE} se
evalúa antes que \texttt{SELECT}).
Sí pueden usarse en \texttt{ORDER BY} y \texttt{HAVING} en MySQL,
aunque esto no es estándar en todos los SGBD.
\end{nota}

%--------------------------------------------------------------
\subsection{\texttt{CONCAT}}
%--------------------------------------------------------------

\texttt{CONCAT} une dos o más cadenas de texto en una sola.
Es una función (no un operador) y acepta cualquier número de
argumentos.

\begin{lstlisting}[caption={CONCAT para construir cadenas compuestas.}]
-- Nombre completo con cargo entre paréntesis
SELECT CONCAT(nombre, ' — ', cargo) AS presentacion
FROM   empleados;

-- CONCAT con separador usando CONCAT_WS
-- (WithSeparator): el primer argumento es el separador
SELECT CONCAT_WS(', ', nombre, departamento, cargo) AS ficha
FROM   empleados;

-- Combinar texto literal con valores de columnas
SELECT CONCAT('El empleado ', nombre,
              ' gana $', FORMAT(salario, 2),
              ' al mes.') AS resumen
FROM   empleados
WHERE  cargo = 'Gerente';
\end{lstlisting}

\begin{nota}
En MySQL, si \textbf{cualquier} argumento de \texttt{CONCAT} es
\texttt{NULL}, el resultado completo es \texttt{NULL}.
Para evitarlo, usa \texttt{COALESCE}: \\
\texttt{CONCAT(COALESCE(nombre, 'Sin nombre'), ' — ', cargo)}.
\texttt{CONCAT\_WS} es más tolerante: ignora los argumentos
\texttt{NULL} en lugar de devolver \texttt{NULL}.
\end{nota}

%--------------------------------------------------------------
\subsection{\texttt{GROUP BY}}
%--------------------------------------------------------------

\texttt{GROUP BY} agrupa las filas que tienen valores idénticos en
las columnas especificadas, para que las funciones de agregación
(\texttt{COUNT}, \texttt{SUM}, \texttt{AVG}, etc.) operen sobre
cada grupo por separado en lugar de sobre toda la tabla.

\begin{definicion}{\texttt{GROUP BY}}
Cláusula que divide las filas del resultado de \texttt{WHERE} en
grupos, donde cada grupo comparte el mismo valor en las columnas
de agrupamiento.
Cada grupo produce exactamente \textbf{una fila} en el resultado.
\end{definicion}

\begin{lstlisting}[caption={GROUP BY con funciones de agregación.}]
-- Número de empleados y salario promedio por departamento
SELECT departamento,
       COUNT(*)           AS num_empleados,
       ROUND(AVG(salario), 2) AS promedio,
       SUM(salario)       AS masa_salarial,
       MIN(salario)       AS minimo,
       MAX(salario)       AS maximo
FROM   empleados
GROUP BY departamento
ORDER BY masa_salarial DESC;
\end{lstlisting}

\begin{lstlisting}[caption={GROUP BY con múltiples columnas.}]
-- Combinaciones únicas de departamento + cargo con conteo
SELECT departamento,
       cargo,
       COUNT(*) AS cantidad
FROM   empleados
GROUP BY departamento, cargo
ORDER BY departamento, cargo;
\end{lstlisting}

\begin{nota}
En SQL estándar (y en MySQL con \texttt{ONLY\_FULL\_GROUP\_BY}
activado), \textbf{toda columna} que aparezca en \texttt{SELECT}
y no esté dentro de una función de agregación debe estar en
\texttt{GROUP BY}.
Si no, el SGBD no sabe qué valor mostrar de entre las múltiples
filas del grupo.
\end{nota}

%--------------------------------------------------------------
\subsection{\texttt{HAVING}}
%--------------------------------------------------------------

\texttt{HAVING} filtra \textbf{grupos} (no filas individuales).
Es el equivalente de \texttt{WHERE} pero aplicado después de que
\texttt{GROUP BY} ha formado los grupos.

\begin{definicion}{\texttt{HAVING}}
Cláusula que filtra los grupos producidos por \texttt{GROUP BY},
conservando sólo aquellos cuya condición (generalmente sobre
una función de agregación) es \texttt{TRUE}.
\end{definicion}

\begin{lstlisting}[caption={HAVING: filtrar grupos por condición agregada.}]
-- Departamentos donde el salario promedio supera 45,000
SELECT departamento,
       ROUND(AVG(salario), 2) AS promedio
FROM   empleados
GROUP BY departamento
HAVING AVG(salario) > 45000;

-- Departamentos con más de 2 empleados
SELECT departamento,
       COUNT(*) AS num_empleados
FROM   empleados
GROUP BY departamento
HAVING COUNT(*) > 2;

-- Combinando WHERE (filtra filas) y HAVING (filtra grupos)
-- Departamentos con más de 1 analista
SELECT departamento,
       COUNT(*) AS num_analistas
FROM   empleados
WHERE  cargo = 'Analista'        -- primero filtra por cargo
GROUP BY departamento
HAVING COUNT(*) > 1;             -- luego filtra grupos
\end{lstlisting}

\begin{ejemplo}
La diferencia entre \texttt{WHERE} y \texttt{HAVING} puede
resumirse así:
\texttt{WHERE} descarta filas \textit{antes} de agrupar
(actúa sobre datos individuales);
\texttt{HAVING} descarta grupos \textit{después} de agrupar
(actúa sobre resultados de agregación).
Intentar usar una función de agregación en \texttt{WHERE} produce
un error de sintaxis.
\end{ejemplo}

%--------------------------------------------------------------
\subsection{\texttt{CASE}}
%--------------------------------------------------------------

\texttt{CASE} es la expresión de \textbf{lógica condicional} de SQL:
permite evaluar condiciones dentro de una consulta y devolver
valores diferentes según cuál se cumpla.
Equivale al \texttt{if/else} o al \texttt{switch} de los lenguajes
de programación.

Existen dos formas sintácticas:

\begin{lstlisting}[caption={CASE simple (compara un valor contra una lista).}]
-- Traduce el nombre del departamento al inglés
SELECT nombre,
       CASE departamento
           WHEN 'Ventas'    THEN 'Sales'
           WHEN 'TI'        THEN 'IT'
           WHEN 'RRHH'      THEN 'HR'
           WHEN 'Finanzas'  THEN 'Finance'
           ELSE 'Unknown'
       END AS department_en
FROM empleados;
\end{lstlisting}

\begin{lstlisting}[caption={CASE buscado (evalúa condiciones booleanas).}]
-- Clasifica el salario en bandas
SELECT nombre,
       salario,
       CASE
           WHEN salario >= 55000 THEN 'Banda A (Senior)'
           WHEN salario >= 45000 THEN 'Banda B (Mid)'
           WHEN salario >= 35000 THEN 'Banda C (Junior)'
           ELSE                       'Banda D (Entrada)'
       END AS banda_salarial
FROM   empleados
ORDER BY salario DESC;
\end{lstlisting}

\begin{lstlisting}[caption={CASE dentro de funciones de agregación.}]
-- Cuenta cuántos empleados hay en cada banda salarial
SELECT
    COUNT(CASE WHEN salario >= 55000 THEN 1 END) AS banda_a,
    COUNT(CASE WHEN salario BETWEEN 45000 AND 54999 THEN 1 END) AS banda_b,
    COUNT(CASE WHEN salario BETWEEN 35000 AND 44999 THEN 1 END) AS banda_c,
    COUNT(CASE WHEN salario < 35000  THEN 1 END) AS banda_d
FROM empleados;
\end{lstlisting}

%--------------------------------------------------------------
\subsection{\texttt{IFNULL} y \texttt{COALESCE}}
%--------------------------------------------------------------

Ambas funciones permiten sustituir un valor \texttt{NULL} por un
valor alternativo.
Son esenciales para producir resultados legibles cuando los datos
pueden ser incompletos.

\begin{lstlisting}[caption={IFNULL: sustituye NULL por un valor por defecto.}]
-- Si telefono es NULL, muestra 'Sin registro'
SELECT nombre,
       IFNULL(telefono, 'Sin registro') AS telefono
FROM   empleados;

-- IFNULL en cálculos: evita que NULL infecte el resultado
SELECT nombre,
       salario + IFNULL(bono, 0) AS salario_total
FROM   empleados;
\end{lstlisting}

\begin{lstlisting}[caption={COALESCE: devuelve el primer valor no NULL de una lista.}]
-- Devuelve el primer contacto disponible entre tres opciones
SELECT nombre,
       COALESCE(telefono_celular,
                telefono_fijo,
                email,
                'Sin contacto') AS contacto_preferido
FROM   empleados;
\end{lstlisting}

\begin{nota}
\texttt{IFNULL(a, b)} es equivalente a \texttt{COALESCE(a, b)},
pero \texttt{COALESCE} acepta cualquier número de argumentos y
es parte del estándar SQL.
Prefiere \texttt{COALESCE} para mayor portabilidad entre SGBD.
\end{nota}

%--------------------------------------------------------------
\subsection{Otros Modificadores}
%--------------------------------------------------------------

\subsubsection{\texttt{DISTINCT} con funciones de agregación}

\texttt{DISTINCT} puede usarse \textit{dentro} de funciones de
agregación para contar o sumar solo valores únicos:

\begin{lstlisting}[caption={DISTINCT dentro de funciones de agregación.}]
-- Cuántos cargos diferentes existen en la empresa
SELECT COUNT(DISTINCT cargo) AS cargos_distintos
FROM   empleados;

-- Suma de salarios únicos (ignora repetidos)
SELECT SUM(DISTINCT salario) AS suma_salarios_unicos
FROM   empleados;
\end{lstlisting}

\subsubsection{\texttt{ROLLUP}: subtotales y total general}

\texttt{WITH ROLLUP} agrega filas extra de subtotales a un
resultado de \texttt{GROUP BY}:

\begin{lstlisting}[caption={GROUP BY con ROLLUP para subtotales.}]
-- Masa salarial por departamento + total general al final
SELECT
    COALESCE(departamento, 'TOTAL EMPRESA') AS nivel,
    COUNT(*)      AS empleados,
    SUM(salario)  AS masa_salarial
FROM   empleados
GROUP BY departamento WITH ROLLUP;
-- La última fila tendrá NULL en departamento (el gran total)
-- COALESCE lo convierte en 'TOTAL EMPRESA'
\end{lstlisting}

\subsubsection{Expresiones aritméticas en \texttt{GROUP BY} y \texttt{HAVING}}

\begin{lstlisting}[caption={Expresiones calculadas combinadas con agrupamiento.}]
-- Año de ingreso como criterio de agrupamiento
SELECT YEAR(fecha_ingreso)   AS anio,
       COUNT(*)              AS contrataciones,
       SUM(salario)          AS masa_salarial_ese_anio
FROM   empleados
GROUP BY YEAR(fecha_ingreso)
HAVING COUNT(*) >= 2
ORDER BY anio DESC;
\end{lstlisting}

\subsubsection{\texttt{FORMAT}: formato de números en la salida}

\begin{lstlisting}[caption={FORMAT para presentación de números.}]
-- Salarios con separador de miles y 2 decimales
SELECT nombre,
       FORMAT(salario, 2)        AS salario_fmt,    -- '48,000.00'
       FORMAT(salario * 12, 0)   AS anual_fmt       -- '576,000'
FROM   empleados
ORDER BY salario DESC;
\end{lstlisting}


% ─── MÓDULO 11: Almacenamiento de datos relacionados ─────
\newpage
\section{Almacenamiento de Datos Relacionados}

Cuando una base de datos posee \textbf{integridad referencial} activa, el orden
y la forma en que se insertan los registros importan tanto como el contenido
mismo. Violar el orden de inserción provoca errores de clave foránea que el
motor rechaza antes de escribir cualquier dato.

\begin{nota}
  Regla general: insertar siempre primero la tabla \textbf{padre} (la que
  contiene la clave primaria referenciada) y después la tabla \textbf{hijo}
  (la que contiene la clave foránea). Para eliminar, el orden es el inverso.
\end{nota}

% ─────────────────────────────────────────────────────────
\subsection{Inserción en Tablas con Relación 1:1}

\begin{definicion}{Relación 1:1}
  Una fila de la tabla \textbf{A} se corresponde con \textbf{exactamente una}
  fila de la tabla \textbf{B}, y viceversa. Se implementa colocando la clave
  foránea en la tabla secundaria con restricción \texttt{UNIQUE}.
\end{definicion}

\textbf{Esquema de ejemplo:}

\begin{lstlisting}[style=sql, caption={Tablas con relación 1:1 — empleado y expediente}]
CREATE TABLE empleados (
    empleado_id   INT          UNSIGNED AUTO_INCREMENT PRIMARY KEY,
    nombre        VARCHAR(100) NOT NULL,
    departamento  VARCHAR(60)  NOT NULL
);

CREATE TABLE expedientes (
    expediente_id INT          UNSIGNED AUTO_INCREMENT PRIMARY KEY,
    empleado_id   INT          UNSIGNED NOT NULL UNIQUE,   -- UNIQUE fuerza 1:1
    nss           CHAR(11)     NOT NULL,
    fecha_alta    DATE         NOT NULL,
    FOREIGN KEY (empleado_id) REFERENCES empleados(empleado_id)
        ON DELETE CASCADE
        ON UPDATE CASCADE
);
\end{lstlisting}

\textbf{Orden correcto de inserción:}

\begin{lstlisting}[style=sql, caption={Inserción respetando integridad referencial 1:1}]
-- Paso 1: insertar en la tabla PADRE
INSERT INTO empleados (nombre, departamento)
VALUES ('Laura Mendoza', 'Ingeniería'),
       ('Roberto Solis', 'Finanzas');

-- Paso 2: recuperar los IDs generados (o usar LAST_INSERT_ID() para uno solo)
-- empleado_id = 1 -> Laura Mendoza
-- empleado_id = 2 -> Roberto Solis

-- Paso 3: insertar en la tabla HIJO referenciando los IDs existentes
INSERT INTO expedientes (empleado_id, nss, fecha_alta)
VALUES (1, '578-12-3456', '2021-03-15'),
       (2, '412-87-9900', '2019-07-01');
\end{lstlisting}

\begin{ejemplo}
  Uso de \texttt{LAST\_INSERT\_ID()} para encadenar inserciones individuales
  de forma segura:

\begin{lstlisting}[style=sql, caption={Encadenamiento con LAST\_INSERT\_ID()}]
INSERT INTO empleados (nombre, departamento)
VALUES ('Carmen Vidal', 'Diseño');

-- LAST_INSERT_ID() devuelve el AUTO_INCREMENT de la última inserción
INSERT INTO expedientes (empleado_id, nss, fecha_alta)
VALUES (LAST_INSERT_ID(), '301-55-7812', '2023-01-10');
\end{lstlisting}
\end{ejemplo}

\begin{nota}
  \texttt{LAST\_INSERT\_ID()} es \textbf{seguro en entornos concurrentes}: cada
  conexión al servidor mantiene su propio contador, por lo que dos sesiones
  simultáneas no interfieren entre sí.
\end{nota}

% ─────────────────────────────────────────────────────────
\subsection{Inserción en Tablas con Relación 1:N}

\begin{definicion}{Relación 1:N}
  Una fila de la tabla \textbf{padre} puede asociarse con \textbf{múltiples}
  filas de la tabla \textbf{hijo}, pero cada fila hijo pertenece a un único
  padre. Es la cardinalidad más frecuente en bases de datos relacionales.
\end{definicion}

\textbf{Esquema de ejemplo:}

\begin{lstlisting}[style=sql, caption={Tablas con relación 1:N — cliente y pedidos}]
CREATE TABLE clientes (
    cliente_id  INT          UNSIGNED AUTO_INCREMENT PRIMARY KEY,
    nombre      VARCHAR(100) NOT NULL,
    correo      VARCHAR(120) NOT NULL UNIQUE
);

CREATE TABLE pedidos (
    pedido_id   INT          UNSIGNED AUTO_INCREMENT PRIMARY KEY,
    cliente_id  INT          UNSIGNED NOT NULL,
    fecha       DATE         NOT NULL,
    total       DECIMAL(10,2) NOT NULL DEFAULT 0.00,
    FOREIGN KEY (cliente_id) REFERENCES clientes(cliente_id)
        ON DELETE RESTRICT
        ON UPDATE CASCADE
);
\end{lstlisting}

\textbf{Inserción de registros dependientes:}

\begin{lstlisting}[style=sql, caption={Inserción 1:N — un cliente con varios pedidos}]
-- Paso 1: insertar el cliente (tabla padre)
INSERT INTO clientes (nombre, correo)
VALUES ('Ana Torres', 'ana.torres@correo.com');

-- Paso 2: insertar sus pedidos (tabla hijo)
--         Todos referencian el mismo cliente_id
INSERT INTO pedidos (cliente_id, fecha, total)
VALUES
    (LAST_INSERT_ID(), '2024-01-15', 1250.00),
    (LAST_INSERT_ID(), '2024-03-02',  340.50),
    (LAST_INSERT_ID(), '2024-06-18', 2890.75);
\end{lstlisting}

\begin{ejemplo}
  Inserción de pedidos para un cliente ya existente, identificado por correo:

\begin{lstlisting}[style=sql, caption={INSERT ... SELECT para referenciar por atributo}]
INSERT INTO pedidos (cliente_id, fecha, total)
SELECT cliente_id, '2024-09-10', 780.00
FROM   clientes
WHERE  correo = 'ana.torres@correo.com';
\end{lstlisting}

  Este patrón evita tener que conocer el \texttt{cliente\_id} numérico de
  antemano y es útil en scripts de carga masiva.
\end{ejemplo}

\begin{nota}
  La política \texttt{ON DELETE RESTRICT} impide eliminar un cliente mientras
  tenga pedidos asociados. Para poder borrarlo habría que eliminar primero sus
  pedidos, o bien cambiar la política a \texttt{ON DELETE CASCADE}.
\end{nota}

% ─────────────────────────────────────────────────────────
\subsection{Inserción en Tablas con Relación N:M}

\begin{definicion}{Relación N:M}
  Múltiples filas de la tabla \textbf{A} pueden relacionarse con múltiples
  filas de la tabla \textbf{B}. Se implementa mediante una
  \textbf{tabla intermedia} (también llamada tabla de unión o pivote) que
  contiene, al menos, las claves foráneas de ambas entidades.
\end{definicion}

\textbf{Esquema de ejemplo:}

\begin{lstlisting}[style=sql, caption={Tablas con relación N:M — estudiantes y cursos}]
CREATE TABLE estudiantes (
    estudiante_id INT          UNSIGNED AUTO_INCREMENT PRIMARY KEY,
    nombre        VARCHAR(100) NOT NULL
);

CREATE TABLE cursos (
    curso_id      INT          UNSIGNED AUTO_INCREMENT PRIMARY KEY,
    titulo        VARCHAR(120) NOT NULL,
    creditos      TINYINT      UNSIGNED NOT NULL
);

-- Tabla intermedia: un estudiante puede tomar N cursos
--                  un curso puede tener N estudiantes
CREATE TABLE inscripciones (
    estudiante_id INT      UNSIGNED NOT NULL,
    curso_id      INT      UNSIGNED NOT NULL,
    fecha_alta    DATE     NOT NULL,
    calificacion  DECIMAL(4,2) DEFAULT NULL,
    PRIMARY KEY (estudiante_id, curso_id),    -- PK compuesta evita duplicados
    FOREIGN KEY (estudiante_id) REFERENCES estudiantes(estudiante_id)
        ON DELETE CASCADE ON UPDATE CASCADE,
    FOREIGN KEY (curso_id)      REFERENCES cursos(curso_id)
        ON DELETE CASCADE ON UPDATE CASCADE
);
\end{lstlisting}

\textbf{Orden de inserción en tres pasos:}

\begin{lstlisting}[style=sql, caption={Paso 1 — poblar ambas tablas padre}]
-- Tabla padre A
INSERT INTO estudiantes (nombre)
VALUES ('Miguel Ríos'),
       ('Sofía Castillo'),
       ('Héctor Núñez');

-- Tabla padre B
INSERT INTO cursos (titulo, creditos)
VALUES ('Bases de Datos',    6),
       ('Álgebra Lineal',    4),
       ('Redes de Cómputo',  5);
\end{lstlisting}

\begin{lstlisting}[style=sql, caption={Paso 2 — insertar en la tabla intermedia}]
-- IDs resultantes (ejemplo):
--   estudiante_id: 1=Miguel, 2=Sofía, 3=Héctor
--   curso_id:      1=Bases de Datos, 2=Álgebra, 3=Redes

INSERT INTO inscripciones (estudiante_id, curso_id, fecha_alta)
VALUES
    (1, 1, '2024-08-20'),   -- Miguel  -> Bases de Datos
    (1, 3, '2024-08-20'),   -- Miguel  -> Redes de Cómputo
    (2, 1, '2024-08-21'),   -- Sofía   -> Bases de Datos
    (2, 2, '2024-08-21'),   -- Sofía   -> Álgebra Lineal
    (3, 1, '2024-08-22'),   -- Héctor  -> Bases de Datos
    (3, 2, '2024-08-22'),   -- Héctor  -> Álgebra Lineal
    (3, 3, '2024-08-22');   -- Héctor  -> Redes de Cómputo
\end{lstlisting}

\begin{lstlisting}[style=sql, caption={Paso 3 — actualizar atributos propios de la relación}]
-- La calificacion es un atributo de la RELACIÓN, no de ninguna entidad sola
UPDATE inscripciones
SET    calificacion = 9.2
WHERE  estudiante_id = 1
  AND  curso_id      = 1;
\end{lstlisting}

\begin{ejemplo}
  Verificar las inscripciones de un estudiante con una consulta de control:

\begin{lstlisting}[style=sql, caption={Consulta de verificación post-inserción N:M}]
SELECT e.nombre        AS estudiante,
       c.titulo        AS curso,
       i.fecha_alta,
       i.calificacion
FROM   inscripciones AS i
INNER JOIN estudiantes AS e ON i.estudiante_id = e.estudiante_id
INNER JOIN cursos      AS c ON i.curso_id      = c.curso_id
WHERE  e.nombre = 'Miguel Ríos';
\end{lstlisting}
\end{ejemplo}

\begin{nota}
  La \textbf{clave primaria compuesta} \texttt{(estudiante\_id, curso\_id)} en
  la tabla intermedia garantiza que un mismo par no pueda insertarse dos veces,
  cumpliendo la regla de integridad de entidad sin necesidad de un campo
  \texttt{AUTO\_INCREMENT} adicional.
\end{nota}

% ── Resumen del módulo ────────────────────────────────────
\subsection*{Resumen del módulo}

\begin{table}[H]
  \centering
  \caption{Estrategias de inserción según la cardinalidad de la relación}
  \renewcommand{\arraystretch}{1.3}
  \rowcolors{2}{gray!10}{white}
  \begin{tabularx}{\textwidth}{l l X}
    \toprule
    \textbf{Relación} & \textbf{Tablas involucradas}     & \textbf{Consideración clave}                              \\
    \midrule
    1:1               & Padre + hijo con FK \texttt{UNIQUE} & Insertar padre primero; usar \texttt{LAST\_INSERT\_ID()}  \\
    1:N               & Padre + uno o más hijos          & Todos los hijos referencian el mismo padre                \\
    N:M               & Dos padres + tabla intermedia    & Insertar ambos padres antes de poblar la tabla pivote     \\
    \bottomrule
  \end{tabularx}
\end{table}
\newpage
\section{Escritura de Datos}

Las sentencias de escritura forman parte del sublenguaje \textbf{DML}
(\textit{Data Manipulation Language}) de SQL. Permiten insertar, modificar y
eliminar filas en las tablas de una base de datos relacional.

\begin{nota}
  Toda operación DML es transaccional. En motores como MySQL/InnoDB o PostgreSQL,
  los cambios pueden deshacerse con \texttt{ROLLBACK} mientras la transacción
  no haya sido confirmada con \texttt{COMMIT}.
\end{nota}

% ─────────────────────────────────────────────────────────
\subsection{\texttt{INSERT}}

\begin{definicion}{INSERT}
  Agrega \textbf{una o más filas} nuevas a una tabla. Los valores proporcionados
  deben respetar los tipos de datos, restricciones \texttt{NOT NULL} y claves
  foráneas definidas en el esquema.
\end{definicion}

\textbf{Sintaxis — una sola fila:}

\begin{lstlisting}[style=sql, caption={INSERT de una fila con columnas explícitas}]
INSERT INTO nombre_tabla (col1, col2, col3)
VALUES (val1, val2, val3);
\end{lstlisting}

\begin{lstlisting}[style=sql, caption={INSERT sin lista de columnas (no recomendado)}]
-- Los valores deben seguir exactamente el orden de la tabla
INSERT INTO nombre_tabla
VALUES (val1, val2, val3);
\end{lstlisting}

\begin{nota}
  Especificar siempre la lista de columnas hace el código más robusto ante futuros
  cambios en el esquema (columnas nuevas, reordenamiento, etc.).
\end{nota}

\textbf{Sintaxis — múltiples filas:}

\begin{lstlisting}[style=sql, caption={INSERT de varias filas en una sola sentencia}]
INSERT INTO nombre_tabla (col1, col2, col3)
VALUES
    (val1a, val2a, val3a),
    (val1b, val2b, val3b),
    (val1c, val2c, val3c);
\end{lstlisting}

\begin{ejemplo}
  Poblar la tabla \texttt{productos} con tres artículos en una sola operación:

\begin{lstlisting}[style=sql, caption={Inserción múltiple en la tabla productos}]
INSERT INTO productos (nombre, precio, stock, categoria_id)
VALUES
    ('Teclado mecánico',  850.00, 40, 3),
    ('Monitor 24"',      3200.00, 15, 3),
    ('Silla ergonómica', 5400.00,  8, 5);
\end{lstlisting}

  Los campos marcados con \texttt{AUTO\_INCREMENT} (\texttt{producto\_id}) se
  omiten; el motor los genera automáticamente.
\end{ejemplo}

\textbf{Inserción a partir de un \texttt{SELECT}:}

\begin{lstlisting}[style=sql, caption={INSERT ... SELECT: copiar filas desde otra tabla}]
INSERT INTO productos_archivo (nombre, precio, categoria_id)
SELECT nombre, precio, categoria_id
FROM   productos
WHERE  stock = 0;
\end{lstlisting}

% ─────────────────────────────────────────────────────────
\subsection{\texttt{UPDATE}}

\begin{definicion}{UPDATE}
  Modifica los valores de \textbf{una o más columnas} en las filas que satisfacen
  la condición \texttt{WHERE}. Si se omite \texttt{WHERE}, la actualización afecta
  a \textbf{todas} las filas de la tabla.
\end{definicion}

\textbf{Sintaxis general:}

\begin{lstlisting}[style=sql, caption={Sintaxis de UPDATE con WHERE}]
UPDATE nombre_tabla
SET    col1 = nuevo_val1,
       col2 = nuevo_val2
WHERE  condicion;
\end{lstlisting}

\begin{ejemplo}
  Aplicar un incremento de 10\,\% al precio de todos los productos de la
  categoría 3:

\begin{lstlisting}[style=sql, caption={UPDATE con expresión aritmética}]
UPDATE productos
SET    precio = precio * 1.10
WHERE  categoria_id = 3;
\end{lstlisting}
\end{ejemplo}

\begin{ejemplo}
  Actualizar varios campos de un registro específico:

\begin{lstlisting}[style=sql, caption={UPDATE de múltiples columnas por clave primaria}]
UPDATE productos
SET    nombre  = 'Teclado mecánico RGB',
       stock   = 55,
       precio  = 920.00
WHERE  producto_id = 1;
\end{lstlisting}
\end{ejemplo}

\begin{nota}
  \textbf{Buenas prácticas antes de ejecutar un UPDATE:}
  \begin{enumerate}
    \item Verificar la condición con un \texttt{SELECT} previo usando el mismo
          \texttt{WHERE}.
    \item Ejecutar dentro de una transacción explícita para poder hacer
          \texttt{ROLLBACK} si el resultado no es el esperado.
    \item En MySQL, activar \texttt{sql\_safe\_updates = 1} para bloquear
          \texttt{UPDATE} sin \texttt{WHERE} o sin clave en la condición.
  \end{enumerate}
\end{nota}

\textbf{UPDATE con JOIN} (MySQL / MariaDB):

\begin{lstlisting}[style=sql, caption={UPDATE que referencia otra tabla mediante JOIN}]
UPDATE productos AS p
INNER JOIN categorias AS c
    ON p.categoria_id = c.categoria_id
SET p.precio = p.precio * 0.90
WHERE c.nombre = 'Electrónica';
\end{lstlisting}

% ─────────────────────────────────────────────────────────
\subsection{\texttt{DELETE}}

\begin{definicion}{DELETE}
  Elimina las filas que cumplen la condición \texttt{WHERE}. Al igual que
  \texttt{UPDATE}, si se omite \texttt{WHERE} se eliminan \textbf{todas} las
  filas de la tabla, aunque la estructura y los metadatos se conservan.
\end{definicion}

\textbf{Sintaxis general:}

\begin{lstlisting}[style=sql, caption={Sintaxis de DELETE con WHERE}]
DELETE FROM nombre_tabla
WHERE condicion;
\end{lstlisting}

\begin{ejemplo}
  Eliminar los productos sin stock que pertenecen a la categoría 5:

\begin{lstlisting}[style=sql, caption={DELETE con condición compuesta}]
DELETE FROM productos
WHERE stock = 0
  AND categoria_id = 5;
\end{lstlisting}
\end{ejemplo}

\textbf{DELETE con JOIN} (MySQL / MariaDB):

\begin{lstlisting}[style=sql, caption={DELETE que filtra mediante una tabla relacionada}]
DELETE p
FROM   productos AS p
INNER JOIN categorias AS c
    ON p.categoria_id = c.categoria_id
WHERE  c.activa = 0;
\end{lstlisting}

% ── DELETE vs TRUNCATE ────────────────────────────────────
\subsubsection*{Diferencias entre \texttt{DELETE} y \texttt{TRUNCATE}}

\begin{table}[H]
  \centering
  \caption{Comparación entre \texttt{DELETE} y \texttt{TRUNCATE}}
  \renewcommand{\arraystretch}{1.3}
  \rowcolors{2}{gray!10}{white}
  \begin{tabularx}{\textwidth}{>{\bfseries}l X X}
    \toprule
    \textbf{Característica}     & \textbf{\texttt{DELETE}}                        & \textbf{\texttt{TRUNCATE}}                   \\
    \midrule
    Sublenguaje SQL             & DML                                             & DDL                                          \\
    Cláusula \texttt{WHERE}     & Sí — elimina filas selectivamente               & No — elimina todas las filas                 \\
    Transaccional               & Sí — admite \texttt{ROLLBACK}                   & Depende del motor; en MySQL no               \\
    Velocidad                   & Más lenta (registro fila a fila)                & Más rápida (libera páginas de datos)         \\
    Reinicio de \texttt{AUTO\_INCREMENT} & No — el contador no se reinicia       & Sí — el contador vuelve a 1                  \\
    Activación de \textit{triggers} & Sí                                         & No en la mayoría de motores                  \\
    Restricciones de FK         & Respeta \texttt{ON DELETE} definido en el FK   & Falla si existen referencias activas         \\
    \bottomrule
  \end{tabularx}
\end{table}

\begin{lstlisting}[style=sql, caption={Diferencia práctica entre DELETE y TRUNCATE}]
-- Elimina solo las filas del año 2020 (reversible)
DELETE FROM ventas
WHERE YEAR(fecha) = 2020;

-- Vacía la tabla completa y reinicia el auto_increment (irreversible en MySQL)
TRUNCATE TABLE ventas;
\end{lstlisting}

\begin{nota}
  Ante la duda entre \texttt{DELETE} y \texttt{TRUNCATE}, preferir
  \texttt{DELETE} dentro de una transacción explícita: ofrece mayor control y
  trazabilidad. Reservar \texttt{TRUNCATE} para entornos de desarrollo o
  procesos de carga masiva donde la velocidad es crítica y no se requiere
  reversibilidad.
\end{nota}

% ── Resumen del módulo ────────────────────────────────────
\subsection*{Resumen del módulo}

\begin{table}[H]
  \centering
  \caption{Sentencias DML de escritura}
  \renewcommand{\arraystretch}{1.3}
  \rowcolors{2}{gray!10}{white}
  \begin{tabularx}{\textwidth}{>{\ttfamily}l l X}
    \toprule
    \textbf{Sentencia}  & \textbf{Acción}          & \textbf{Consideración clave}                          \\
    \midrule
    INSERT              & Agregar filas            & Especificar siempre la lista de columnas              \\
    INSERT \ldots SELECT & Copiar filas            & Útil para archivado o migración de datos              \\
    UPDATE              & Modificar filas          & Nunca omitir \texttt{WHERE}; verificar con SELECT previo \\
    DELETE              & Eliminar filas           & Transaccional; respeta triggers y claves foráneas     \\
    TRUNCATE            & Vaciar tabla completa    & DDL, más rápido, pero no reversible en MySQL          \\
    \bottomrule
  \end{tabularx}
\end{table}



% ─── MÓDULO 7: Administración de la base de datos ────────
\newpage
\section{Administración de la Base de Datos}

Las sentencias de administración de bases de datos pertenecen al sublenguaje
\textbf{DDL}. Operan sobre el contenedor de más alto nivel en MySQL: la
\textbf{base de datos} (también llamada \textit{schema}), que agrupa tablas,
vistas, procedimientos y demás objetos.

% ─────────────────────────────────────────────────────────
\subsection{\texttt{CREATE DATABASE}}

\begin{definicion}{CREATE DATABASE}
  Crea una nueva base de datos vacía en el servidor. Permite especificar el
  \textbf{juego de caracteres} (\textit{character set}) y la
  \textbf{colación} (\textit{collation}) que determinarán cómo se almacenan
  y comparan las cadenas de texto en todas sus tablas.
\end{definicion}

\begin{lstlisting}[style=sql, caption={Sintaxis de CREATE DATABASE}]
CREATE DATABASE [IF NOT EXISTS] nombre_bd
    [CHARACTER SET  charset_name]
    [COLLATE        collation_name];
\end{lstlisting}

\begin{ejemplo}
  Crear una base de datos para una aplicación en español con soporte
  completo de Unicode:

\begin{lstlisting}[style=sql, caption={Base de datos con utf8mb4 y collation en español}]
CREATE DATABASE IF NOT EXISTS tienda_online
    CHARACTER SET  utf8mb4
    COLLATE        utf8mb4_unicode_ci;

-- Seleccionar la base de datos activa
USE tienda_online;
\end{lstlisting}
\end{ejemplo}

\textbf{Juegos de caracteres y colaciones más comunes:}

\begin{table}[H]
  \centering
  \caption{Character sets y collations habituales en MySQL}
  \renewcommand{\arraystretch}{1.3}
  \rowcolors{2}{gray!10}{white}
  \begin{tabularx}{\textwidth}{>{\ttfamily}l >{\ttfamily}l X}
    \toprule
    \textbf{charset}  & \textbf{collation}          & \textbf{Uso recomendado}                        \\
    \midrule
    utf8mb4  & utf8mb4\_unicode\_ci  & Unicode completo; soporta emojis; recomendado   \\
    utf8mb4  & utf8mb4\_0900\_ai\_ci & Unicode 9.0; sin acento/mayúsculas; MySQL 8.0+  \\
    utf8mb4  & utf8mb4\_bin          & Comparación binaria; distingue mayúsculas       \\
    latin1   & latin1\_swedish\_ci   & Legado; solo caracteres ISO-8859-1              \\
    \bottomrule
  \end{tabularx}
\end{table}

\begin{nota}
  Usar \texttt{utf8} (sin \texttt{mb4}) en MySQL es un error frecuente: ese
  alias solo soporta hasta 3 bytes por carácter y excluye emojis y algunos
  ideogramas CJK. La elección correcta para proyectos nuevos es siempre
  \texttt{utf8mb4}.
\end{nota}

\begin{lstlisting}[style=sql, caption={Comandos de exploración de bases de datos}]
SHOW DATABASES;                    -- listar todas las bases de datos
SHOW CREATE DATABASE tienda_online; -- ver la definición completa
SELECT DATABASE();                  -- ver la base de datos activa
\end{lstlisting}

% ─────────────────────────────────────────────────────────
\subsection{\texttt{DROP DATABASE}}

\begin{definicion}{DROP DATABASE}
  Elimina la base de datos y \textbf{todos los objetos que contiene}
  (tablas, vistas, procedimientos, datos) de forma permanente e
  irreversible. No existe confirmación ni papelera de reciclaje.
\end{definicion}

\begin{lstlisting}[style=sql, caption={Sintaxis de DROP DATABASE}]
DROP DATABASE [IF EXISTS] nombre_bd;
\end{lstlisting}

\begin{ejemplo}
\begin{lstlisting}[style=sql, caption={DROP DATABASE con y sin guarda}]
-- Genera error si la base de datos no existe
DROP DATABASE tienda_online;

-- No genera error si no existe (útil en scripts de despliegue)
DROP DATABASE IF EXISTS tienda_online;
\end{lstlisting}
\end{ejemplo}

\begin{nota}
  \textbf{Precauciones antes de ejecutar DROP DATABASE:}
  \begin{enumerate}
    \item Verificar que se tiene un respaldo (\textit{backup}) actualizado.
    \item Confirmar que ninguna aplicación activa tiene conexiones abiertas
          a esa base de datos.
    \item Comprobar el nombre con \texttt{SHOW DATABASES} antes de ejecutar.
    \item En producción, revocar el privilegio \texttt{DROP} a usuarios de
          aplicación para evitar ejecuciones accidentales.
  \end{enumerate}
\end{nota}

% ─── MÓDULO 13: Conceptos avanzados ──────────────────────
\newpage
\section{Conceptos Avanzados}

Los objetos descritos en este módulo permiten construir bases de datos más
eficientes, seguras y mantenibles. Aunque no son imprescindibles para
consultas básicas, son herramientas fundamentales en entornos de producción.

% ─────────────────────────────────────────────────────────
\subsection{\texttt{INDEX}}

\begin{definicion}{Índice}
  Estructura de datos auxiliar que el motor mantiene sincronizada con la
  tabla para acelerar las operaciones de búsqueda, ordenamiento y unión.
  Mejora el rendimiento de \texttt{SELECT} a cambio de un coste adicional
  en \texttt{INSERT}, \texttt{UPDATE} y \texttt{DELETE}, y de espacio en
  disco.
\end{definicion}

\textbf{Tipos de índice en MySQL:}

\begin{table}[H]
  \centering
  \caption{Tipos de índice disponibles en MySQL/InnoDB}
  \renewcommand{\arraystretch}{1.3}
  \rowcolors{2}{gray!10}{white}
  \begin{tabularx}{\textwidth}{>{\ttfamily}l l X}
    \toprule
    \textbf{Tipo}   & \textbf{Motor}       & \textbf{Uso recomendado}                              \\
    \midrule
    B-TREE          & InnoDB, MyISAM       & Búsquedas exactas, rangos, \texttt{ORDER BY}          \\
    HASH            & Memory               & Solo igualdad exacta; no soporta rangos               \\
    FULLTEXT        & InnoDB (5.6+)        & Búsqueda de texto libre en columnas \texttt{TEXT}     \\
    SPATIAL         & InnoDB, MyISAM       & Datos geoespaciales (\texttt{GEOMETRY})               \\
    \bottomrule
  \end{tabularx}
\end{table}

\begin{lstlisting}[style=sql, caption={Creación y eliminación de índices}]
-- Índice simple en una columna
CREATE INDEX idx_apellido ON clientes(apellido);

-- Índice compuesto (el orden de columnas importa)
CREATE INDEX idx_cat_precio ON productos(categoria_id, precio);

-- Índice UNIQUE (equivale a restricción UNIQUE)
CREATE UNIQUE INDEX uq_correo ON usuarios(correo);

-- Índice FULLTEXT para búsqueda de texto
CREATE FULLTEXT INDEX ft_descripcion ON productos(descripcion);

-- Eliminar un índice
DROP INDEX idx_apellido ON clientes;
\end{lstlisting}

\begin{lstlisting}[style=sql, caption={Uso del índice FULLTEXT con MATCH ... AGAINST}]
SELECT nombre, descripcion
FROM   productos
WHERE  MATCH(descripcion) AGAINST('teclado mecánico' IN BOOLEAN MODE);
\end{lstlisting}

\begin{nota}
  \textbf{Cuándo crear un índice:}
  \begin{itemize}
    \item Columnas que aparecen frecuentemente en \texttt{WHERE}, \texttt{JOIN ON},
          \texttt{ORDER BY} o \texttt{GROUP BY}.
    \item Columnas con alta cardinalidad (muchos valores distintos).
  \end{itemize}
  \textbf{Cuándo evitarlo:}
  \begin{itemize}
    \item Tablas pequeñas (el escaneo completo es más rápido).
    \item Columnas con muy baja cardinalidad (\texttt{activo}, \texttt{genero}).
    \item Tablas con escrituras masivas y pocas lecturas.
  \end{itemize}
\end{nota}

% ─────────────────────────────────────────────────────────
\subsection{\texttt{TRIGGER}}

\begin{definicion}{TRIGGER}
  Bloque de código SQL que se ejecuta \textbf{automáticamente} en respuesta
  a un evento DML (\texttt{INSERT}, \texttt{UPDATE} o \texttt{DELETE}) sobre
  una tabla, antes (\texttt{BEFORE}) o después (\texttt{AFTER}) de que el
  cambio se aplique.
\end{definicion}

\begin{lstlisting}[style=sql, caption={Sintaxis general de un TRIGGER}]
CREATE TRIGGER nombre_trigger
    { BEFORE | AFTER } { INSERT | UPDATE | DELETE }
    ON nombre_tabla
    FOR EACH ROW
BEGIN
    -- cuerpo del trigger
    -- NEW.columna  -> valor nuevo (INSERT / UPDATE)
    -- OLD.columna  -> valor anterior (UPDATE / DELETE)
END;
\end{lstlisting}

\begin{ejemplo}
  Registrar en una tabla de auditoría cada cambio de precio en
  \texttt{productos}:

\begin{lstlisting}[style=sql, caption={TRIGGER de auditoría de cambios de precio}]
CREATE TABLE auditoria_precios (
    id          INT UNSIGNED AUTO_INCREMENT PRIMARY KEY,
    producto_id INT UNSIGNED NOT NULL,
    precio_ant  DECIMAL(10,2),
    precio_new  DECIMAL(10,2),
    cambiado_en DATETIME NOT NULL DEFAULT CURRENT_TIMESTAMP
);

DELIMITER $$
CREATE TRIGGER trg_auditoria_precio
    AFTER UPDATE ON productos
    FOR EACH ROW
BEGIN
    IF OLD.precio <> NEW.precio THEN
        INSERT INTO auditoria_precios
               (producto_id, precio_ant, precio_new)
        VALUES (NEW.producto_id, OLD.precio, NEW.precio);
    END IF;
END$$
DELIMITER ;
\end{lstlisting}
\end{ejemplo}

\begin{nota}
  Los triggers dificultan la depuración y pueden crear efectos secundarios
  inesperados si no se documentan. Úsarlos con moderación; para lógica
  compleja, preferir procedimientos almacenados invocados explícitamente.
\end{nota}

% ─────────────────────────────────────────────────────────
\subsection{\texttt{VIEW}}

\begin{definicion}{Vista (\textit{View})}
  Consulta \texttt{SELECT} almacenada con nombre propio que puede usarse
  como si fuera una tabla. No duplica datos: cada vez que se consulta, el
  motor ejecuta la consulta subyacente. Simplifica consultas complejas,
  encapsula lógica y puede restringir el acceso a columnas sensibles.
\end{definicion}

\begin{lstlisting}[style=sql, caption={Creación y consulta de una vista}]
-- Crear la vista
CREATE VIEW vista_pedidos_activos AS
    SELECT p.pedido_id,
           c.nombre          AS cliente,
           p.total,
           p.fecha
    FROM   pedidos    AS p
    INNER JOIN clientes AS c ON p.cliente_id = c.cliente_id
    WHERE  p.estado = 'activo';

-- Usar la vista exactamente como una tabla
SELECT cliente, SUM(total) AS total_compras
FROM   vista_pedidos_activos
GROUP  BY cliente
ORDER  BY total_compras DESC;

-- Reemplazar o eliminar la vista
CREATE OR REPLACE VIEW vista_pedidos_activos AS ...;
DROP VIEW IF EXISTS vista_pedidos_activos;
\end{lstlisting}

\begin{nota}
  Las vistas en MySQL no almacenan datos en caché por defecto. Para
  escenarios de alto rendimiento con consultas agregadas costosas,
  considerar \textbf{vistas materializadas} (disponibles en PostgreSQL) o
  tablas de resumen actualizadas por eventos programados.
\end{nota}

% ─────────────────────────────────────────────────────────
\subsection{\texttt{STORED PROCEDURE}}

\begin{definicion}{Procedimiento almacenado}
  Bloque de código SQL con nombre, almacenado en el servidor y ejecutable
  mediante \texttt{CALL}. Acepta parámetros de entrada (\texttt{IN}),
  salida (\texttt{OUT}) o ambos (\texttt{INOUT}), y permite estructuras de
  control (\texttt{IF}, \texttt{WHILE}, \texttt{LOOP}, cursores).
\end{definicion}

\begin{lstlisting}[style=sql, caption={Procedimiento con parámetros IN y OUT}]
DELIMITER $$
CREATE PROCEDURE sp_resumen_cliente(
    IN  p_cliente_id INT UNSIGNED,
    OUT p_total_pedidos INT,
    OUT p_monto_total   DECIMAL(12,2)
)
BEGIN
    SELECT COUNT(*),
           COALESCE(SUM(total), 0)
    INTO   p_total_pedidos,
           p_monto_total
    FROM   pedidos
    WHERE  cliente_id = p_cliente_id;
END$$
DELIMITER ;

-- Invocar el procedimiento
CALL sp_resumen_cliente(7, @num_pedidos, @monto);
SELECT @num_pedidos AS pedidos, @monto AS monto_total;
\end{lstlisting}

\begin{ejemplo}
  Procedimiento con flujo de control para aplicar descuentos escalonados:

\begin{lstlisting}[style=sql, caption={Stored procedure con IF / ELSEIF}]
DELIMITER $$
CREATE PROCEDURE sp_aplicar_descuento(
    IN p_producto_id INT UNSIGNED,
    IN p_porcentaje  DECIMAL(5,2)
)
BEGIN
    IF p_porcentaje <= 0 OR p_porcentaje > 80 THEN
        SIGNAL SQLSTATE '45000'
            SET MESSAGE_TEXT = 'Porcentaje fuera de rango permitido';
    ELSEIF p_porcentaje <= 20 THEN
        UPDATE productos
        SET    precio = precio * (1 - p_porcentaje / 100)
        WHERE  producto_id = p_producto_id;
    ELSE
        -- Descuento mayor requiere aprobación: solo registrar solicitud
        INSERT INTO solicitudes_descuento(producto_id, porcentaje)
        VALUES (p_producto_id, p_porcentaje);
    END IF;
END$$
DELIMITER ;
\end{lstlisting}
\end{ejemplo}

% ─────────────────────────────────────────────────────────
\subsection{Transacciones}

\begin{definicion}{Transacción}
  Unidad lógica de trabajo que agrupa una o más operaciones DML. Garantiza
  las propiedades \textbf{ACID}: Atomicidad, Consistencia, Aislamiento y
  Durabilidad. Si alguna operación falla, todo el grupo puede deshacerse.
\end{definicion}

\begin{table}[H]
  \centering
  \caption{Comandos de control de transacciones}
  \renewcommand{\arraystretch}{1.3}
  \rowcolors{2}{gray!10}{white}
  \begin{tabularx}{\textwidth}{>{\ttfamily}l X}
    \toprule
    \textbf{Comando}          & \textbf{Efecto}                                               \\
    \midrule
    START TRANSACTION / BEGIN & Inicia una transacción explícita                              \\
    COMMIT                    & Confirma y persiste todos los cambios                         \\
    ROLLBACK                  & Deshace todos los cambios desde el último COMMIT              \\
    SAVEPOINT nombre          & Crea un punto de restauración parcial dentro de la transacción\\
    ROLLBACK TO nombre        & Deshace hasta el savepoint indicado (sin terminar la transacción)\\
    RELEASE SAVEPOINT nombre  & Elimina un savepoint sin hacer rollback                       \\
    \bottomrule
  \end{tabularx}
\end{table}

\begin{lstlisting}[style=sql, caption={Transacción con SAVEPOINT y manejo de errores}]
START TRANSACTION;

    INSERT INTO pedidos (cliente_id, total, fecha)
    VALUES (12, 1500.00, CURDATE());

    SAVEPOINT sp_pedido_insertado;

    INSERT INTO detalle_pedido (pedido_id, producto_id, cantidad)
    VALUES (LAST_INSERT_ID(), 5, 3);

    -- Si este UPDATE falla, se puede revertir solo el detalle
    UPDATE productos SET stock = stock - 3 WHERE producto_id = 5;

    -- Todo correcto: confirmar
    COMMIT;

-- En caso de error detectado en la aplicación:
-- ROLLBACK TO sp_pedido_insertado;  -- deshace solo desde el savepoint
-- ROLLBACK;                          -- deshace toda la transacción
\end{lstlisting}

\begin{nota}
  MySQL opera en modo \texttt{AUTOCOMMIT = 1} por defecto: cada sentencia
  DML es una transacción implícita que se confirma de inmediato.
  \texttt{START TRANSACTION} suspende ese modo hasta que se ejecute
  \texttt{COMMIT} o \texttt{ROLLBACK}.
\end{nota}

% ─────────────────────────────────────────────────────────
\subsection{Concurrencia}

\begin{definicion}{Concurrencia}
  Capacidad del motor para atender \textbf{múltiples transacciones
  simultáneas} sin que interfieran entre sí. Se gestiona mediante
  \textbf{bloqueos} (\textit{locks}) y \textbf{niveles de aislamiento}
  (\textit{isolation levels}).
\end{definicion}

\textbf{Niveles de aislamiento en SQL estándar:}

\begin{table}[H]
  \centering
  \caption{Niveles de aislamiento y fenómenos que previenen}
  \renewcommand{\arraystretch}{1.3}
  \rowcolors{2}{gray!10}{white}
  \begin{tabularx}{\textwidth}{l X c c c}
    \toprule
    \textbf{Nivel}      & \textbf{Descripción resumida}
      & \textbf{\textit{Dirty read}}
      & \textbf{\textit{Non-rep. read}}
      & \textbf{\textit{Phantom}} \\
    \midrule
    READ UNCOMMITTED & Ve cambios no confirmados de otras transacciones   & Sí & Sí & Sí \\
    READ COMMITTED   & Solo ve cambios ya confirmados                     & No & Sí & Sí \\
    REPEATABLE READ  & Garantiza lecturas consistentes (defecto InnoDB)   & No & No & Sí \\
    SERIALIZABLE     & Ejecución secuencial total; máximo aislamiento      & No & No & No \\
    \bottomrule
  \end{tabularx}
\end{table}

\begin{lstlisting}[style=sql, caption={Configurar y consultar el nivel de aislamiento}]
-- Ver el nivel actual
SELECT @@transaction_isolation;

-- Cambiar para la sesión actual
SET SESSION TRANSACTION ISOLATION LEVEL REPEATABLE READ;

-- Cambiar globalmente (requiere privilegio SUPER)
SET GLOBAL  TRANSACTION ISOLATION LEVEL READ COMMITTED;
\end{lstlisting}

\begin{nota}
  \textbf{Deadlock:} dos transacciones se bloquean mutuamente esperando
  recursos que la otra retiene. InnoDB detecta el ciclo automáticamente y
  aborta la transacción de menor coste. La aplicación debe estar preparada
  para reintentar la operación al recibir el error \texttt{ER\_LOCK\_DEADLOCK}
  (\texttt{1213}). Para minimizarlos: acceder siempre a las tablas en el
  mismo orden y mantener las transacciones lo más cortas posible.
\end{nota}

% ─── MÓDULO 14: Conexión desde código ────────────────────
\newpage
\section{Conexión desde Código}

Integrar una base de datos en una aplicación requiere un \textbf{conector}
o \textit{driver}: biblioteca que traduce las llamadas del lenguaje de
programación al protocolo de red del motor SQL.

% ─────────────────────────────────────────────────────────
\subsection{Conectores y Drivers}

\begin{table}[H]
  \centering
  \caption{Conectores oficiales y populares por lenguaje}
  \renewcommand{\arraystretch}{1.3}
  \rowcolors{2}{gray!10}{white}
  \begin{tabularx}{\textwidth}{l >{\ttfamily}l X}
    \toprule
    \textbf{Lenguaje} & \textbf{Librería / Driver}  & \textbf{Instalación / Referencia}          \\
    \midrule
    Python   & mysql-connector-python & \texttt{pip install mysql-connector-python} \\
    Python   & PyMySQL                & \texttt{pip install pymysql}                \\
    Python   & SQLAlchemy (ORM)       & \texttt{pip install sqlalchemy}             \\
    Node.js  & mysql2                 & \texttt{npm install mysql2}                 \\
    Node.js  & Prisma (ORM)           & \texttt{npm install prisma}                 \\
    Java     & MySQL Connector/J      & Dependencia Maven: \texttt{mysql:mysql-connector-java} \\
    PHP      & PDO\_MySQL             & Incluido en PHP; activar extensión \texttt{pdo\_mysql} \\
    Go       & go-sql-driver/mysql    & \texttt{go get github.com/go-sql-driver/mysql} \\
    \bottomrule
  \end{tabularx}
\end{table}

\begin{ejemplo}
  Conexión y consulta básica con Python (\texttt{mysql-connector-python}):

\begin{lstlisting}[language=Python, basicstyle=\ttfamily\small,
    keywordstyle=\bfseries\color{NavyBlue},
    commentstyle=\itshape\color{Gray},
    backgroundcolor=\color{gray!6}, frame=single,
    caption={Conexión a MySQL desde Python}]
import mysql.connector

conn = mysql.connector.connect(
    host     = "localhost",
    user     = "app_user",
    password = "s3cr3to",
    database = "tienda_online"
)

cursor = conn.cursor(dictionary=True)
cursor.execute("SELECT producto_id, nombre, precio FROM productos LIMIT 5")

for fila in cursor.fetchall():
    print(fila["nombre"], fila["precio"])

cursor.close()
conn.close()
\end{lstlisting}
\end{ejemplo}

\begin{ejemplo}
  Conexión con Node.js usando \texttt{mysql2} y \texttt{async/await}:

\begin{lstlisting}[language=Java, basicstyle=\ttfamily\small,
    keywordstyle=\bfseries\color{NavyBlue},
    commentstyle=\itshape\color{Gray},
    backgroundcolor=\color{gray!6}, frame=single,
    caption={Conexión a MySQL desde Node.js}]
const mysql = require('mysql2/promise');

const conn = await mysql.createConnection({
    host:     'localhost',
    user:     'app_user',
    password: 's3cr3to',
    database: 'tienda_online'
});

const [rows] = await conn.execute(
    'SELECT nombre, precio FROM productos WHERE categoria_id = ?',
    [3]   // parámetro enlazado (previene SQL Injection)
);

console.table(rows);
await conn.end();
\end{lstlisting}
\end{ejemplo}

% ─────────────────────────────────────────────────────────
\subsection{SQL Injection}

\begin{definicion}{SQL Injection}
  Vulnerabilidad que permite a un atacante \textbf{inyectar código SQL
  arbitrario} en una consulta, manipulando su lógica para extraer,
  modificar o destruir datos, o para eludir la autenticación.
\end{definicion}

\begin{ejemplo}
  Código vulnerable: concatenación directa de entrada del usuario.

\begin{lstlisting}[language=Python, basicstyle=\ttfamily\small,
    keywordstyle=\bfseries\color{NavyBlue},
    commentstyle=\itshape\color{Gray},
    backgroundcolor=\color{gray!6}, frame=single,
    caption={Código VULNERABLE a SQL Injection}]
usuario = input("Usuario: ")   # atacante ingresa:  ' OR '1'='1
query = "SELECT * FROM usuarios WHERE username = '" + usuario + "'"
# Consulta resultante:
# SELECT * FROM usuarios WHERE username = '' OR '1'='1'
# -> devuelve TODOS los usuarios sin contraseña
cursor.execute(query)
\end{lstlisting}
\end{ejemplo}

\begin{ejemplo}
  Código seguro: consultas parametrizadas (\textit{prepared statements}).

\begin{lstlisting}[language=Python, basicstyle=\ttfamily\small,
    keywordstyle=\bfseries\color{NavyBlue},
    commentstyle=\itshape\color{Gray},
    backgroundcolor=\color{gray!6}, frame=single,
    caption={Prevención con consultas parametrizadas}]
usuario = input("Usuario: ")
# El driver escapa automáticamente el valor
cursor.execute(
    "SELECT * FROM usuarios WHERE username = %s",
    (usuario,)   # parámetro enlazado, nunca concatenado
)
\end{lstlisting}
\end{ejemplo}

\textbf{Medidas de prevención complementarias:}

\begin{itemize}
  \item Usar \textbf{siempre} consultas parametrizadas o ORMs que las
        generen internamente.
  \item Aplicar el \textbf{principio de mínimo privilegio}: el usuario de
        la aplicación solo debe tener los permisos estrictamente necesarios
        (\texttt{SELECT}, \texttt{INSERT}, etc.), nunca \texttt{DROP} o
        \texttt{GRANT}.
  \item Validar y sanitizar la entrada en la capa de aplicación antes de
        pasarla al driver.
  \item Evitar exponer mensajes de error detallados del motor en respuestas
        de la API.
\end{itemize}

% ─── MÓDULO 15: Otros clientes y sistemas ────────────────
\newpage
\section{Otros Clientes SQL}

Existen múltiples herramientas gráficas y de línea de comandos para
interactuar con bases de datos SQL sin necesidad de escribir directamente
en la terminal.

\begin{table}[H]
  \centering
  \caption{Clientes SQL populares}
  \renewcommand{\arraystretch}{1.3}
  \rowcolors{2}{gray!10}{white}
  \begin{tabularx}{\textwidth}{l l X}
    \toprule
    \textbf{Cliente}  & \textbf{Plataforma}          & \textbf{Características destacadas}                  \\
    \midrule
    DBeaver           & Windows, macOS, Linux        & Gratuito y de código abierto; soporta más de 80 motores \\
    TablePlus         & macOS, Windows, Linux, iOS   & Interfaz nativa moderna; licencia de pago con prueba gratuita \\
    DataGrip          & Windows, macOS, Linux        & IDE de JetBrains; refactorización SQL, inspecciones inteligentes \\
    pgAdmin           & Web / escritorio             & Cliente oficial de PostgreSQL; gestión completa del servidor \\
    Adminer           & Web (un solo archivo PHP)    & Ligero, sin instalación; útil en servidores remotos  \\
    MySQL Workbench   & Windows, macOS, Linux        & Cliente oficial de MySQL; diseño ER visual incluido   \\
    \bottomrule
  \end{tabularx}
\end{table}

% ─────────────────────────────────────────────────────────
\section{PostgreSQL}

PostgreSQL es un sistema de gestión de bases de datos relacionales
\textbf{de código abierto}, con más de 35 años de desarrollo activo,
reconocido por su robustez, extensibilidad y cumplimiento estricto del
estándar SQL.

\textbf{Diferencias clave con MySQL:}

\begin{table}[H]
  \centering
  \caption{MySQL vs.\ PostgreSQL — diferencias principales}
  \renewcommand{\arraystretch}{1.3}
  \rowcolors{2}{gray!10}{white}
  \begin{tabularx}{\textwidth}{l X X}
    \toprule
    \textbf{Aspecto}          & \textbf{MySQL}                          & \textbf{PostgreSQL}                        \\
    \midrule
    Licencia                  & GPL / comercial (Oracle)                & PostgreSQL License (MIT-like)              \\
    Modelo MVCC               & Solo en InnoDB                          & Nativo en todo el motor                    \\
    Tipos de datos avanzados  & Limitados                               & JSON/JSONB, Arrays, HSTORE, UUID nativo    \\
    Vistas materializadas     & No nativo                               & Sí, con \texttt{REFRESH MATERIALIZED VIEW} \\
    Full-text search          & Básico                                  & Avanzado; soporta diccionarios y ranking   \\
    Extensiones               & Plugins limitados                       & PostGIS, pgcrypto, TimescaleDB, etc.       \\
    Procedimientos almacenados& SQL / MySQL dialect                     & PL/pgSQL, PL/Python, PL/Perl, etc.        \\
    Secuencias                & \texttt{AUTO\_INCREMENT}                & \texttt{SERIAL} / \texttt{GENERATED ALWAYS AS IDENTITY} \\
    \bottomrule
  \end{tabularx}
\end{table}

\begin{nota}
  PostgreSQL es generalmente preferido para aplicaciones que requieren
  consultas complejas, integridad de datos estricta o tipos de datos
  avanzados. MySQL sigue siendo muy popular en aplicaciones web de alto
  tráfico gracias a su velocidad en lecturas simples y su amplio soporte
  en servicios de alojamiento.
\end{nota}

% ─────────────────────────────────────────────────────────
\section{Despliegue en la Nube}

Los servicios gestionados (\textit{Database as a Service}, DBaaS) eliminan
la carga operativa de instalar, parchear y respaldar el motor de base de
datos.

\begin{table}[H]
  \centering
  \caption{Servicios de base de datos en la nube}
  \renewcommand{\arraystretch}{1.3}
  \rowcolors{2}{gray!10}{white}
  \begin{tabularx}{\textwidth}{l l X}
    \toprule
    \textbf{Servicio}     & \textbf{Motor(es)}           & \textbf{Características destacadas}                     \\
    \midrule
    AWS RDS               & MySQL, PostgreSQL, MariaDB, Oracle, SQL Server
                                                          & Backups automáticos, Multi-AZ, réplicas de lectura      \\
    AWS Aurora            & MySQL / PostgreSQL compatible& Rendimiento hasta 5$\times$ sobre RDS; serverless disponible \\
    Google Cloud SQL      & MySQL, PostgreSQL, SQL Server& Integración con BigQuery y Vertex AI                    \\
    Azure Database        & MySQL, PostgreSQL            & Alta disponibilidad con zona de disponibilidad flexible  \\
    PlanetScale           & MySQL (Vitess)               & Branching de esquemas; escala horizontal automática      \\
    Railway               & MySQL, PostgreSQL, MongoDB   & Despliegue en segundos; ideal para prototipos            \\
    Supabase              & PostgreSQL                   & BaaS open-source; autenticación y API REST incluidos     \\
    Neon                  & PostgreSQL                   & Serverless; branching por pull-request                   \\
    \bottomrule
  \end{tabularx}
\end{table}

% ─────────────────────────────────────────────────────────
\section{Próximos Pasos}

\subsection*{Caminos de especialización}

\begin{itemize}
  \item \textbf{DBA (\textit{Database Administrator}):} administración de
        servidores, alta disponibilidad, replicación, \textit{tuning} de
        rendimiento, backups y recuperación ante desastres.
  \item \textbf{Data Engineer:} pipelines de datos (ETL/ELT), almacenes de
        datos (\textit{data warehouses}), herramientas como dbt, Apache
        Spark, Airflow y servicios como BigQuery o Redshift.
  \item \textbf{Backend Developer:} integración de bases de datos en APIs
        REST y GraphQL, ORM avanzados, patrones de diseño de esquemas,
        caché con Redis.
  \item \textbf{Data Analyst / BI:} consultas analíticas avanzadas,
        funciones de ventana, herramientas de visualización (Metabase,
        Tableau, Power BI).
\end{itemize}

\subsection*{Recursos recomendados}

\begin{itemize}
  \item \textbf{Documentación oficial:}
        \href{https://dev.mysql.com/doc/}{dev.mysql.com/doc} y
        \href{https://www.postgresql.org/docs/}{postgresql.org/docs}.
  \item \textbf{Práctica interactiva:}
        \href{https://sqlzoo.net}{SQLZoo},
        \href{https://www.hackerrank.com/domains/sql}{HackerRank SQL} y
        \href{https://leetcode.com/problemset/database/}{LeetCode Database}.
  \item \textbf{Libros de referencia:}
        \textit{Learning MySQL} (Seyed M.M. \& Jay Pipes),
        \textit{PostgreSQL: Up and Running} (Regina O. Obe \& Leo S. Hsu),
        \textit{Designing Data-Intensive Applications} (Martin Kleppmann).
  \item \textbf{Modelado de datos:}
        \href{https://dbdiagram.io}{dbdiagram.io} para diseño visual de
        esquemas con notación DBML.
\end{itemize}

\printbibliography[heading=bibintoc, title={Referencias}]
\end{document}


